\documentclass[8pt]{beamer}
\usetheme{Madrid}
\usecolortheme{default}
\usepackage{amsmath, amssymb, graphicx, array, multirow, booktabs, xcolor, tikz}
\usetikzlibrary{shapes,arrows,positioning,fit,backgrounds}

% Compact font settings
\setbeamerfont{title}{size=\Large}
\setbeamerfont{subtitle}{size=\small}
\setbeamerfont{author}{size=\scriptsize}
\setbeamerfont{date}{size=\scriptsize}
\setbeamerfont{frametitle}{size=\large}
\setbeamerfont{framesubtitle}{size=\normalsize}
\setbeamerfont{enumerate item}{size=\footnotesize}
\setbeamerfont{itemize item}{size=\footnotesize}
\setbeamerfont{block title}{size=\footnotesize}

% Compact spacing
\setlength{\itemsep}{0.08ex}
\setlength{\parskip}{0.08ex}
\setlength{\leftmargini}{0.3cm}
\setbeamersize{text margin left=3mm, text margin right=3mm}
\addtobeamertemplate{frame begin}{\vspace{-0.4cm}}{}

% Colors
\definecolor{myblue}{RGB}{0,0,255}
\definecolor{myred}{RGB}{255,0,0}
\definecolor{mygreen}{RGB}{0,128,0}
\definecolor{myorange}{RGB}{255,165,0}
\definecolor{mypurple}{RGB}{128,0,128}
\definecolor{mygray}{RGB}{128,128,128}
\definecolor{mygold}{RGB}{255,215,0}

\title{Risk and Risk Factors}
\subtitle{100 Questions: Algorithmic Bias, Fairness, Explainability, GenAI Risks}
\author{GARP RAI Exam Preparation}
\date{}

\begin{document}
	
	\begin{frame}
		\titlepage
	\end{frame}
	
	% ============================================================
	% SECTION 1: INTRODUCTION TO AI RISKS - QUESTIONS 1-5
	% ============================================================
	
	\begin{frame}{Q1: AI Risks Overview - Tutorial}
		\fontsize{6.5}{7.5}\selectfont
		\textbf{QUESTION: Why is understanding AI risks particularly challenging?}
		
		\vspace{0.1cm}
		\textbf{OPTIONS:}
		\begin{enumerate}
			\item AI systems are always perfectly reliable
			\item AI is hugely pervasive in range and variety, with rapidly developing capabilities
			\item AI risks are all well-understood and documented
			\item AI systems never fail
		\end{enumerate}
		
		\vspace{0.1cm}
		\textcolor{myblue}{\textbf{ANSWER: B}}
		
		\vspace{0.1cm}
		\textbf{DETAILED EXPLANATION:}
		\begin{itemize}
			\item AI is pervasive across many domains (finance, healthcare, justice)
			\item AI capabilities are rapidly developing
			\item High-stakes applications mean failures have significant consequences
			\item Risk measurement and mitigation are challenging
			\item Regulatory responses are often slow to develop
		\end{itemize}
		
		\vspace{0.1cm}
		\textcolor{myred}{\textbf{WHY OTHER OPTIONS ARE WRONG:}}
		\begin{itemize}
			\item \textcolor{myred}{A:} AI systems are NOT always perfectly reliable
			\item \textcolor{myred}{C:} Many AI risks are still being discovered
			\item \textcolor{myred}{D:} AI systems can and do fail
		\end{itemize}
		
		\textcolor{myred}{\textbf{EXAM TRAP:}} AI risks are \textcolor{myblue}{evolving, pervasive, and high-stakes}!
	\end{frame}
	
	\begin{frame}{Q2: Pervasiveness Challenge - Tutorial}
		\fontsize{6.5}{7.5}\selectfont
		\textbf{QUESTION: Why does the pervasiveness of AI pose a risk management challenge?}
		
		\vspace{0.1cm}
		\textbf{OPTIONS:}
		\begin{enumerate}
			\item AI is only used in one industry
			\item Risks take different forms depending on the use case
			\item All AI risks are identical
			\item Pervasiveness makes risks easier to predict
		\end{enumerate}
		
		\vspace{0.1cm}
		\textcolor{myblue}{\textbf{ANSWER: B}}
		
		\vspace{0.1cm}
		\textbf{DETAILED EXPLANATION:}
		\begin{itemize}
			\item AI applications span many domains (credit scoring, healthcare, policing)
			\item Each domain has unique risk profiles
			\item A loan denial algorithm has different risks than a medical diagnostic algorithm
			\item Risk prediction must account for context
		\end{itemize}
		
		\vspace{0.1cm}
		\textcolor{myred}{\textbf{WHY OTHER OPTIONS ARE WRONG:}}
		\begin{itemize}
			\item \textcolor{myred}{A:} AI is used across MANY industries
			\item \textcolor{myred}{C:} Risks vary significantly by use case
			\item \textcolor{myred}{D:} Pervasiveness makes risks HARDER to predict
		\end{itemize}
		
		\textcolor{myred}{\textbf{EXAM TRAP:}} AI risks are \textcolor{myblue}{context-dependent} across different use cases!
	\end{frame}
	
	\begin{frame}{Q3: High-Stakes AI Use - Tutorial}
		\fontsize{6.5}{7.5}\selectfont
		\textbf{QUESTION: What is a consequence of deploying AI in high-stakes settings?}
		
		\vspace{0.1cm}
		\textbf{OPTIONS:}
		\begin{enumerate}
			\item Technological shortcomings have minimal impact
			\item Technological shortcomings can have significant consequences for people
			\item AI is never used in high-stakes settings
			\item High-stakes use eliminates all risks
		\end{enumerate}
		
		\vspace{0.1cm}
		\textcolor{myblue}{\textbf{ANSWER: B}}
		
		\vspace{0.1cm}
		\textbf{DETAILED EXPLANATION:}
		\begin{itemize}
			\item Loan applications affect financial well-being
			\item Medical diagnoses affect health outcomes
			\item Bail decisions affect personal freedom
			\item Algorithmic errors can have profound impacts
		\end{itemize}
		
		\vspace{0.1cm}
		\textcolor{myred}{\textbf{WHY OTHER OPTIONS ARE WRONG:}}
		\begin{itemize}
			\item \textcolor{myred}{A:} High-stakes means impacts are SIGNIFICANT
			\item \textcolor{myred}{C:} AI IS used in high-stakes settings
			\item \textcolor{myred}{D:} High-stakes use increases, not eliminates, risk
		\end{itemize}
		
		\textcolor{myred}{\textbf{EXAM TRAP:}} High-stakes AI = \textcolor{myblue}{significant consequences} for individuals!
	\end{frame}
	
	\begin{frame}{Q4: Rapid AI Development - Tutorial}
		\fontsize{6.5}{7.5}\selectfont
		\textbf{QUESTION: What challenge does rapid AI development create?}
		
		\vspace{0.1cm}
		\textbf{OPTIONS:}
		\begin{enumerate}
			\item It makes risk measurement easier
			\item It creates challenges for impact prediction, risk measurement, and governance
			\item It eliminates all regulatory concerns
			\item It slows down innovation
		\end{enumerate}
		
		\vspace{0.1cm}
		\textcolor{myblue}{\textbf{ANSWER: B}}
		
		\vspace{0.1cm}
		\textbf{DETAILED EXPLANATION:}
		\begin{itemize}
			\item AI capabilities advance faster than our understanding of risks
			\item Regulatory responses lag behind technological change
			\item Governance mechanisms are difficult to establish quickly
			\item Slow regulation allows harmful technologies to persist
		\end{itemize}
		
		\vspace{0.1cm}
		\textcolor{myred}{\textbf{WHY OTHER OPTIONS ARE WRONG:}}
		\begin{itemize}
			\item \textcolor{myred}{A:} Rapid development makes measurement HARDER
			\item \textcolor{myred}{C:} Regulatory concerns are significant
			\item \textcolor{myred}{D:} Rapid development ACCELERATES innovation
		\end{itemize}
		
		\textcolor{myred}{\textbf{EXAM TRAP:}} Rapid AI development = \textcolor{myblue}{governance and regulation challenges}!
	\end{frame}
	
	\begin{frame}{Q5: Interconnected AI Risks - Tutorial}
		\fontsize{6.5}{7.5}\selectfont
		\textbf{QUESTION: How are AI risks typically related?}
		
		\vspace{0.1cm}
		\textbf{OPTIONS:}
		\begin{enumerate}
			\item They are completely independent of each other
			\item They are often interconnected - lack of transparency affects fairness, bias leads to reputational risks
			\item There is only one type of AI risk
			\item Risks are never interconnected
		\end{enumerate}
		
		\vspace{0.1cm}
		\textcolor{myblue}{\textbf{ANSWER: B}}
		
		\vspace{0.1cm}
		\textbf{DETAILED EXPLANATION:}
		\begin{itemize}
			\item Transparency issues affect fairness assessment
			\item Algorithmic bias leads to reputational damage
			\item Safety failures can lead to fairness concerns
			\item Risks form a complex network
		\end{itemize}
		
		\vspace{0.1cm}
		\textcolor{myred}{\textbf{WHY OTHER OPTIONS ARE WRONG:}}
		\begin{itemize}
			\item \textcolor{myred}{A:} Risks are highly interconnected
			\item \textcolor{myred}{C:} There are many types of AI risks
			\item \textcolor{myred}{D:} Risks ARE interconnected
		\end{itemize}
		
		\textcolor{myred}{\textbf{EXAM TRAP:}} AI risks are \textcolor{myblue}{interconnected and mutually reinforcing}!
	\end{frame}
	
	% ============================================================
	% SECTION 2: ALGORITHMIC BIAS AND FAIRNESS - QUESTIONS 6-30
	% ============================================================
	
	\begin{frame}{Q6: Algorithmic Bias Definition - Tutorial}
		\fontsize{6.5}{7.5}\selectfont
		\textbf{QUESTION: What is algorithmic bias?}
		
		\vspace{0.1cm}
		\textbf{OPTIONS:}
		\begin{enumerate}
			\item Random errors in algorithmic predictions
			\item A systematic deviation of an algorithm's output, performance, or impact relative to some norm or standard
			\item Algorithms that are always perfectly accurate
			\item Errors that only occur in supervised learning
		\end{enumerate}
		
		\vspace{0.1cm}
		\textcolor{myblue}{\textbf{ANSWER: B}}
		
		\vspace{0.1cm}
		\textbf{DETAILED EXPLANATION:}
		\begin{itemize}
			\item Systematic (not random) deviation
			\item Can affect output, performance, or impact
			\item Relative to some norm, aim, standard, or baseline
			\item Can be explicit (intentional) or implicit (unintended)
		\end{itemize}
		
		\vspace{0.1cm}
		\textcolor{myred}{\textbf{WHY OTHER OPTIONS ARE WRONG:}}
		\begin{itemize}
			\item \textcolor{myred}{A:} Bias is SYSTEMATIC, not random
			\item \textcolor{myred}{C:} Algorithms CAN be biased
			\item \textcolor{myred}{D:} Bias can occur in ANY type of learning
		\end{itemize}
		
		\textcolor{myred}{\textbf{EXAM TRAP:}} Algorithmic bias = \textcolor{myblue}{systematic deviation from a norm}!
	\end{frame}
	
	\begin{frame}{Q7: Explicit vs Implicit Bias - Tutorial}
		\fontsize{6.5}{7.5}\selectfont
		\textbf{QUESTION: What is the difference between explicit and implicit bias?}
		
		\vspace{0.1cm}
		\textbf{OPTIONS:}
		\begin{enumerate}
			\item Explicit bias is always worse than implicit bias
			\item Explicit bias is intentional/projected by developers; implicit bias creeps in undetected
			\item Implicit bias is intentional; explicit bias is accidental
			\item There is no difference
		\end{enumerate}
		
		\vspace{0.1cm}
		\textcolor{myblue}{\textbf{ANSWER: B}}
		
		\vspace{0.1cm}
		\textbf{DETAILED EXPLANATION:}
		\begin{itemize}
			\item \textbf{Explicit bias:} Developer intentionally projects own biases
			\item \textbf{Implicit bias:} Bias creeps in undetected through design or data
			\item Both types can lead to unfair outcomes
			\item Implicit bias is often harder to detect and correct
		\end{itemize}
		
		\vspace{0.1cm}
		\textcolor{myred}{\textbf{WHY OTHER OPTIONS ARE WRONG:}}
		\begin{itemize}
			\item \textcolor{myred}{A:} Both can be equally harmful
			\item \textcolor{myred}{C:} This is reversed
			\item \textcolor{myred}{D:} They differ in intentionality
		\end{itemize}
		
		\textcolor{myred}{\textbf{EXAM TRAP:}} Explicit = \textcolor{myblue}{intentional}; Implicit = \textcolor{myblue}{unintended}!
	\end{frame}
	
	\begin{frame}{Q8: Individual Fairness Definition - Tutorial}
		\fontsize{6.5}{7.5}\selectfont
		\textbf{QUESTION: What is individual fairness in algorithmic design?}
		
		\vspace{0.1cm}
		\textbf{OPTIONS:}
		\begin{enumerate}
			\item Treating all groups equally
			\item The doctrine of "treating like cases alike"
			\item Ensuring equal outcomes across demographics
			\item Always using protected characteristics
		\end{enumerate}
		
		\vspace{0.1cm}
		\textcolor{myblue}{\textbf{ANSWER: B}}
		
		\vspace{0.1cm}
		\textbf{DETAILED EXPLANATION:}
		\begin{itemize}
			\item Based on Aristotelian doctrine
			\item Similarly situated individuals should be treated equally
			\item Focuses on individuals, not groups
			\item Requires agreement on what counts as "alike"
		\end{itemize}
		
		\vspace{0.1cm}
		\textcolor{myred}{\textbf{WHY OTHER OPTIONS ARE WRONG:}}
		\begin{itemize}
			\item \textcolor{myred}{A:} This is group fairness
			\item \textcolor{myred}{C:} This is demographic parity
			\item \textcolor{myred}{D:} Individual fairness avoids using protected characteristics
		\end{itemize}
		
		\textcolor{myred}{\textbf{EXAM TRAP:}} Individual fairness = \textcolor{myblue}{treating like cases alike}!
	\end{frame}
	
	\begin{frame}{Q9: Group Fairness Definition - Tutorial}
		\fontsize{6.5}{7.5}\selectfont
		\textbf{QUESTION: What is group fairness?}
		
		\vspace{0.1cm}
		\textbf{OPTIONS:}
		\begin{enumerate}
			\item Treating each individual uniquely
			\item Examining statistical differences across groups
			\item Only considering individuals with protected characteristics
			\item Ignoring group-level patterns
		\end{enumerate}
		
		\vspace{0.1cm}
		\textcolor{myblue}{\textbf{ANSWER: B}}
		
		\vspace{0.1cm}
		\textbf{DETAILED EXPLANATION:}
		\begin{itemize}
			\item Looks at groups sharing protected characteristics
			\item Asks whether groups are disadvantaged
			\item Examples: admission rates by gender, loan approvals by race
			\item Does not consider individuals directly
		\end{itemize}
		
		\vspace{0.1cm}
		\textcolor{myred}{\textbf{WHY OTHER OPTIONS ARE WRONG:}}
		\begin{itemize}
			\item \textcolor{myred}{A:} This is individual fairness
			\item \textcolor{myred}{C:} Group fairness considers ALL groups
			\item \textcolor{myred}{D:} Group fairness EXAMINES group patterns
		\end{itemize}
		
		\textcolor{myred}{\textbf{EXAM TRAP:}} Group fairness = \textcolor{myblue}{statistical differences across groups}!
	\end{frame}
	
	\begin{frame}{Q10: College Admission Example - Tutorial}
		\fontsize{6.5}{7.5}\selectfont
		\textbf{QUESTION: In the college admission example, what fairness issue arises?}
		
		\vspace{0.1cm}
		\textbf{OPTIONS:}
		\begin{enumerate}
			\item All applicants are treated fairly
			\item Individual fairness requires agreeing on what counts as "alike"
			\item Socioeconomic factors never matter
			\item Algorithms always make perfect decisions
		\end{enumerate}
		
		\vspace{0.1cm}
		\textcolor{myblue}{\textbf{ANSWER: B}}
		
		\vspace{0.1cm}
		\textbf{DETAILED EXPLANATION:}
		\begin{itemize}
			\item What counts as "alike" is often contentious
			\item Should we consider hardships overcome?
			\item How to translate resilience into numbers?
			\item Different people have different views on relevant criteria
		\end{itemize}
		
		\vspace{0.1cm}
		\textcolor{myred}{\textbf{WHY OTHER OPTIONS ARE WRONG:}}
		\begin{itemize}
			\item \textcolor{myred}{A:} Fairness is not automatically achieved
			\item \textcolor{myred}{C:} Socioeconomic factors are often considered
			\item \textcolor{myred}{D:} Algorithms can make unfair decisions
		\end{itemize}
		
		\textcolor{myred}{\textbf{EXAM TRAP:}} Individual fairness requires \textcolor{myblue}{agreement on what is "alike"}!
	\end{frame}
	
	\begin{frame}{Q11: Demographic Parity Definition - Tutorial}
		\fontsize{6.5}{7.5}\selectfont
		\textbf{QUESTION: What is demographic parity?}
		
		\vspace{0.1cm}
		\textbf{OPTIONS:}
		\begin{enumerate}
			\item Predictions must be equally accurate across groups
			\item The distribution of predictions must be identical across subpopulations
			\item Positive cases must be correctly identified equally across groups
			\item Precision must be equal across groups
		\end{enumerate}
		
		\vspace{0.1cm}
		\textcolor{myblue}{\textbf{ANSWER: B}}
		
		\vspace{0.1cm}
		\textbf{DETAILED EXPLANATION:}
		\begin{itemize}
			\item $P(\hat{Y}=1|A=0) = P(\hat{Y}=1|A=1)$
			\item Outcome distribution is identical across groups
			\item Example: same admission rate across genders
			\item Independent of true values (doesn't consider correctness)
		\end{itemize}
		
		\vspace{0.1cm}
		\textcolor{myred}{\textbf{WHY OTHER OPTIONS ARE WRONG:}}
		\begin{itemize}
			\item \textcolor{myred}{A:} This is predictive rate parity/equal opportunity
			\item \textcolor{myred}{C:} This is equal opportunity
			\item \textcolor{myred}{D:} This is predictive rate parity
		\end{itemize}
		
		\textcolor{myred}{\textbf{EXAM TRAP:}} Demographic parity = \textcolor{myblue}{equal distribution of predictions}!
	\end{frame}
	
	\begin{frame}{Q12: Demographic Parity Formula - Tutorial}
		\fontsize{6.5}{7.5}\selectfont
		\textbf{QUESTION: What is the mathematical definition of demographic parity?}
		
		\vspace{0.1cm}
		\textbf{OPTIONS:}
		\begin{enumerate}
			\item $P(Y=1|\hat{Y}=1,A=0) = P(Y=1|\hat{Y}=1,A=1)$
			\item $P(\hat{Y}=1|A=0) = P(\hat{Y}=1|A=1)$
			\item $P(\hat{Y}=1|Y=1,A=0) = P(\hat{Y}=1|Y=1,A=1)$
			\item $P(Y=1|A=0) = P(Y=1|A=1)$
		\end{enumerate}
		
		\vspace{0.1cm}
		\textcolor{myblue}{\textbf{ANSWER: B}}
		
		\vspace{0.1cm}
		\textbf{DETAILED EXPLANATION:}
		\begin{itemize}
			\item $\hat{Y}$ = predicted outcome
			\item $A$ = protected attribute/group membership
			\item Probability of positive prediction given group membership
			\item Does not depend on actual outcomes ($Y$)
		\end{itemize}
		
		\vspace{0.1cm}
		\textcolor{myred}{\textbf{WHY OTHER OPTIONS ARE WRONG:}}
		\begin{itemize}
			\item \textcolor{myred}{A:} This is predictive rate parity (uses precision)
			\item \textcolor{myred}{C:} This is equal opportunity (uses sensitivity)
			\item \textcolor{myred}{D:} This is base rate equality
		\end{itemize}
		
		\textcolor{myred}{\textbf{EXAM TRAP:}} Demographic parity = \textcolor{myblue}{$P(\hat{Y}=1|A=0) = P(\hat{Y}=1|A=1)$}!
	\end{frame}
	
	\begin{frame}{Q13: Demographic Parity Drawback - Tutorial}
		\fontsize{6.5}{7.5}\selectfont
		\textbf{QUESTION: What is a major drawback of demographic parity?}
		
		\vspace{0.1cm}
		\textbf{OPTIONS:}
		\begin{enumerate}
			\item It is always impossible to achieve
			\item It does not consider the quality/accuracy of predictions
			\item It is too computationally expensive
			\item It requires protected characteristics
		\end{enumerate}
		
		\vspace{0.1cm}
		\textcolor{myblue}{\textbf{ANSWER: B}}
		
		\vspace{0.1cm}
		\textbf{DETAILED EXPLANATION:}
		\begin{itemize}
			\item Independent of true values
			\item Can be satisfied even if all predictions are wrong
			\item Example: Unqualified candidates labeled as qualified equally across groups
			\item Doesn't ensure predictions are correct
		\end{itemize}
		
		\vspace{0.1cm}
		\textcolor{myred}{\textbf{WHY OTHER OPTIONS ARE WRONG:}}
		\begin{itemize}
			\item \textcolor{myred}{A:} Not always impossible, but can be challenging
			\item \textcolor{myred}{C:} It's computationally simple
			\item \textcolor{myred}{D:} It does require knowing group membership
		\end{itemize}
		
		\textcolor{myred}{\textbf{EXAM TRAP:}} Demographic parity ignores \textcolor{myblue}{prediction accuracy/quality}!
	\end{frame}
	
	\begin{frame}{Q14: Predictive Rate Parity Definition - Tutorial}
		\fontsize{6.5}{7.5}\selectfont
		\textbf{QUESTION: What is predictive rate parity?}
		
		\vspace{0.1cm}
		\textbf{OPTIONS:}
		\begin{enumerate}
			\item The distribution of predictions is identical across groups
			\item Precision (TP/(TP+FP)) is equal across groups
			\item Sensitivity (TP/(TP+FN)) is equal across groups
			\item Accuracy is equal across groups
		\end{enumerate}
		
		\vspace{0.1cm}
		\textcolor{myblue}{\textbf{ANSWER: B}}
		
		\vspace{0.1cm}
		\textbf{DETAILED EXPLANATION:}
		\begin{itemize}
			\item $P(Y=1|\hat{Y}=1,A=0) = P(Y=1|\hat{Y}=1,A=1)$
			\item Among those predicted positive, proportion correctly identified is same
			\item Also called "positive predictive value parity"
			\item Ensures false positive rates are balanced across groups
		\end{itemize}
		
		\vspace{0.1cm}
		\textcolor{myred}{\textbf{WHY OTHER OPTIONS ARE WRONG:}}
		\begin{itemize}
			\item \textcolor{myred}{A:} This is demographic parity
			\item \textcolor{myred}{C:} This is equal opportunity
			\item \textcolor{myred}{D:} This is accuracy parity (different concept)
		\end{itemize}
		
		\textcolor{myred}{\textbf{EXAM TRAP:}} Predictive rate parity = \textcolor{myblue}{equal precision across groups}!
	\end{frame}
	
	\begin{frame}{Q15: Predictive Rate Parity Drawback - Tutorial}
		\fontsize{6.5}{7.5}\selectfont
		\textbf{QUESTION: What is a drawback of predictive rate parity?}
		
		\vspace{0.1cm}
		\textbf{OPTIONS:}
		\begin{enumerate}
			\item It ignores true positives
			\item It does not account for differences in base rates across groups
			\item It is always impossible to achieve
			\item It requires no data
		\end{enumerate}
		
		\vspace{0.1cm}
		\textcolor{myblue}{\textbf{ANSWER: B}}
		
		\vspace{0.1cm}
		\textbf{DETAILED EXPLANATION:}
		\begin{itemize}
			\item Base rates may differ significantly across groups
			\item Achieving parity becomes extremely challenging
			\item Can have adverse effects on algorithmic performance
			\item May require sacrificing accuracy for fairness
		\end{itemize}
		
		\vspace{0.1cm}
		\textcolor{myred}{\textbf{WHY OTHER OPTIONS ARE WRONG:}}
		\begin{itemize}
			\item \textcolor{myred}{A:} It specifically considers true positives (precision)
			\item \textcolor{myred}{C:} Not always impossible, but can be difficult
			\item \textcolor{myred}{D:} Requires data to calculate
		\end{itemize}
		
		\textcolor{myred}{\textbf{EXAM TRAP:}} Predictive rate parity struggles with \textcolor{myblue}{different base rates}!
	\end{frame}
	
	\begin{frame}{Q16: Equal Opportunity Definition - Tutorial}
		\fontsize{6.5}{7.5}\selectfont
		\textbf{QUESTION: What is equal opportunity?}
		
		\vspace{0.1cm}
		\textbf{OPTIONS:}
		\begin{enumerate}
			\item The distribution of predictions is identical across groups
			\item Precision is equal across groups
			\item Sensitivity (TP/(TP+FN)) is equal across groups
			\item All predictions are always correct
		\end{enumerate}
		
		\vspace{0.1cm}
		\textcolor{myblue}{\textbf{ANSWER: C}}
		
		\vspace{0.1cm}
		\textbf{DETAILED EXPLANATION:}
		\begin{itemize}
			\item $P(\hat{Y}=1|Y=1,A=0) = P(\hat{Y}=1|Y=1,A=1)$
			\item True positive rate is equal across groups
			\item Positive cases equally likely to be identified correctly
			\item Important in medical diagnostics (don't miss disease)
		\end{itemize}
		
		\vspace{0.1cm}
		\textcolor{myred}{\textbf{WHY OTHER OPTIONS ARE WRONG:}}
		\begin{itemize}
			\item \textcolor{myred}{A:} This is demographic parity
			\item \textcolor{myred}{B:} This is predictive rate parity
			\item \textcolor{myred}{D:} No model is always correct
		\end{itemize}
		
		\textcolor{myred}{\textbf{EXAM TRAP:}} Equal opportunity = \textcolor{myblue}{equal true positive rate (sensitivity)}!
	\end{frame}
	
	\begin{frame}{Q17: Equal Opportunity Formula - Tutorial}
		\fontsize{6.5}{7.5}\selectfont
		\textbf{QUESTION: What is the mathematical definition of equal opportunity?}
		
		\vspace{0.1cm}
		\textbf{OPTIONS:}
		\begin{enumerate}
			\item $P(\hat{Y}=1|A=0) = P(\hat{Y}=1|A=1)$
			\item $P(Y=1|\hat{Y}=1,A=0) = P(Y=1|\hat{Y}=1,A=1)$
			\item $P(\hat{Y}=1|Y=1,A=0) = P(\hat{Y}=1|Y=1,A=1)$
			\item $P(Y=1|A=0) = P(Y=1|A=1)$
		\end{enumerate}
		
		\vspace{0.1cm}
		\textcolor{myblue}{\textbf{ANSWER: C}}
		
		\vspace{0.1cm}
		\textbf{DETAILED EXPLANATION:}
		\begin{itemize}
			\item $\hat{Y}$ = predicted outcome
			\item $Y$ = actual outcome
			\item $A$ = protected attribute
			\item Probability of correct positive prediction given actual positive
		\end{itemize}
		
		\vspace{0.1cm}
		\textcolor{myred}{\textbf{WHY OTHER OPTIONS ARE WRONG:}}
		\begin{itemize}
			\item \textcolor{myred}{A:} This is demographic parity
			\item \textcolor{myred}{B:} This is predictive rate parity
			\item \textcolor{myred}{D:} This is base rate equality
		\end{itemize}
		
		\textcolor{myred}{\textbf{EXAM TRAP:}} Equal opportunity = \textcolor{myblue}{$P(\hat{Y}=1|Y=1,A=0) = P(\hat{Y}=1|Y=1,A=1)$}!
	\end{frame}
	
	\begin{frame}{Q18: Equalized Odds Definition - Tutorial}
		\fontsize{6.5}{7.5}\selectfont
		\textbf{QUESTION: What is equalized odds?}
		
		\vspace{0.1cm}
		\textbf{OPTIONS:}
		\begin{enumerate}
			\item Equal distribution of predictions
			\item Equal precision across groups
			\item Equal true positive rate AND equal false positive rate across groups
			\item Equal accuracy across groups
		\end{enumerate}
		
		\vspace{0.1cm}
		\textcolor{myblue}{\textbf{ANSWER: C}}
		
		\vspace{0.1cm}
		\textbf{DETAILED EXPLANATION:}
		\begin{itemize}
			\item More restrictive than equal opportunity
			\item Requires both sensitivity AND specificity equality
			\item $P(\hat{Y}=1|Y=1,A=0) = P(\hat{Y}=1|Y=1,A=1)$
			\item $P(\hat{Y}=1|Y=0,A=0) = P(\hat{Y}=1|Y=0,A=1)$
		\end{itemize}
		
		\vspace{0.1cm}
		\textcolor{myred}{\textbf{WHY OTHER OPTIONS ARE WRONG:}}
		\begin{itemize}
			\item \textcolor{myred}{A:} This is demographic parity
			\item \textcolor{myred}{B:} This is predictive rate parity
			\item \textcolor{myred}{D:} Accuracy parity is different
		\end{itemize}
		
		\textcolor{myred}{\textbf{EXAM TRAP:}} Equalized odds = \textcolor{myblue}{equal TPR AND equal FPR}!
	\end{frame}
	
	\begin{frame}{Q19: Fairness Impossibility Theorem - Tutorial}
		\fontsize{6.5}{7.5}\selectfont
		\textbf{QUESTION: What does the fairness impossibility theorem state?}
		
		\vspace{0.1cm}
		\textbf{OPTIONS:}
		\begin{enumerate}
			\item All fairness measures can be satisfied simultaneously
			\item It is impossible to satisfy demographic parity, predictive rate parity, and equal opportunity simultaneously when base rates differ
			\item Fairness is always achievable
			\item Only one fairness measure exists
		\end{enumerate}
		
		\vspace{0.1cm}
		\textcolor{myblue}{\textbf{ANSWER: B}}
		
		\vspace{0.1cm}
		\textbf{DETAILED EXPLANATION:}
		\begin{itemize}
			\item When base rates differ between groups
			\item Cannot satisfy all three fairness notions
			\item Trade-offs are unavoidable
			\item Must deliberate on which fairness measure to prioritize
		\end{itemize}
		
		\vspace{0.1cm}
		\textcolor{myred}{\textbf{WHY OTHER OPTIONS ARE WRONG:}}
		\begin{itemize}
			\item \textcolor{myred}{A:} This is impossible (the theorem proves it)
			\item \textcolor{myred}{C:} Fairness has inherent trade-offs
			\item \textcolor{myred}{D:} There are multiple fairness measures
		\end{itemize}
		
		\textcolor{myred}{\textbf{EXAM TRAP:}} Cannot satisfy \textcolor{myblue}{all fairness measures} simultaneously when base rates differ!
	\end{frame}
	
	\begin{frame}{Q20: Accuracy Limitation - Tutorial}
		\fontsize{6.5}{7.5}\selectfont
		\textbf{QUESTION: Why is accuracy not always a good performance measure?}
		
		\vspace{0.1cm}
		\textbf{OPTIONS:}
		\begin{enumerate}
			\item It is always a perfect measure
			\item It can be misleading with unbalanced data
			\item It is too difficult to calculate
			\item It only works for regression
		\end{enumerate}
		
		\vspace{0.1cm}
		\textcolor{myblue}{\textbf{ANSWER: B}}
		
		\vspace{0.1cm}
		\textbf{DETAILED EXPLANATION:}
		\begin{itemize}
			\item Example: 100 applicants, 90 unsuitable, 10 suitable
			\item Algorithm classifies all as unsuitable: Accuracy = 90\%
			\item But algorithm provides no useful information (all unsuitable)
			\item Especially problematic in unbalanced datasets
		\end{itemize}
		
		\vspace{0.1cm}
		\textcolor{myred}{\textbf{WHY OTHER OPTIONS ARE WRONG:}}
		\begin{itemize}
			\item \textcolor{myred}{A:} Accuracy has significant limitations
			\item \textcolor{myred}{C:} Accuracy is easy to calculate
			\item \textcolor{myred}{D:} Accuracy is for classification, not regression
		\end{itemize}
		
		\textcolor{myred}{\textbf{EXAM TRAP:}} Accuracy is \textcolor{myblue}{misleading with unbalanced data}!
	\end{frame}
	
	\begin{frame}{Q21: Precision Definition - Tutorial}
		\fontsize{6.5}{7.5}\selectfont
		\textbf{QUESTION: What is precision in classification?}
		
		\vspace{0.1cm}
		\textbf{OPTIONS:}
		\begin{enumerate}
			\item TP/(TP+FN)
			\item TP/(TP+FP)
			\item (TP+TN)/(TP+TN+FP+FN)
			\item TN/(TN+FP)
		\end{enumerate}
		
		\vspace{0.1cm}
		\textcolor{myblue}{\textbf{ANSWER: B}}
		
		\vspace{0.1cm}
		\textbf{DETAILED EXPLANATION:}
		\begin{itemize}
			\item Proportion of positive predictions that are correct
			\item Also called Positive Predictive Value (PPV)
			\item $Precision = \frac{TP}{TP+FP}$
			\item Answers: "Of those predicted positive, how many were actually positive?"
		\end{itemize}
		
		\vspace{0.1cm}
		\textcolor{myred}{\textbf{WHY OTHER OPTIONS ARE WRONG:}}
		\begin{itemize}
			\item \textcolor{myred}{A:} This is recall/sensitivity
			\item \textcolor{myred}{C:} This is accuracy
			\item \textcolor{myred}{D:} This is specificity
		\end{itemize}
		
		\textcolor{myred}{\textbf{EXAM TRAP:}} Precision = \textcolor{myblue}{TP/(TP+FP)} - "of predicted positives, how many were right?"
	\end{frame}
	
	\begin{frame}{Q22: Sensitivity/Recall Definition - Tutorial}
		\fontsize{6.5}{7.5}\selectfont
		\textbf{QUESTION: What is sensitivity (recall) in classification?}
		
		\vspace{0.1cm}
		\textbf{OPTIONS:}
		\begin{enumerate}
			\item TP/(TP+FN)
			\item TP/(TP+FP)
			\item (TP+TN)/(TP+TN+FP+FN)
			\item TN/(TN+FP)
		\end{enumerate}
		
		\vspace{0.1cm}
		\textcolor{myblue}{\textbf{ANSWER: A}}
		
		\vspace{0.1cm}
		\textbf{DETAILED EXPLANATION:}
		\begin{itemize}
			\item Proportion of actual positives that are correctly identified
			\item Also called True Positive Rate (TPR)
			\item $Sensitivity = \frac{TP}{TP+FN}$
			\item Answers: "Of those actually positive, how many did we catch?"
		\end{itemize}
		
		\vspace{0.1cm}
		\textcolor{myred}{\textbf{WHY OTHER OPTIONS ARE WRONG:}}
		\begin{itemize}
			\item \textcolor{myred}{B:} This is precision
			\item \textcolor{myred}{C:} This is accuracy
			\item \textcolor{myred}{D:} This is specificity
		\end{itemize}
		
		\textcolor{myred}{\textbf{EXAM TRAP:}} Sensitivity = \textcolor{myblue}{TP/(TP+FN)} - "of actual positives, how many did we catch?"
	\end{frame}
	
	\begin{frame}{Q23: Precision vs Sensitivity Trade-off - Tutorial}
		\fontsize{6.5}{7.5}\selectfont
		\textbf{QUESTION: What trade-off exists between precision and sensitivity?}
		
		\vspace{0.1cm}
		\textbf{OPTIONS:}
		\begin{enumerate}
			\item They always move in the same direction
			\item There is often a trade-off; maximizing one can reduce the other
			\item They are completely independent
			\item Precision is always more important
		\end{enumerate}
		
		\vspace{0.1cm}
		\textcolor{myblue}{\textbf{ANSWER: B}}
		
		\vspace{0.1cm}
		\textbf{DETAILED EXPLANATION:}
		\begin{itemize}
			\item To catch all positives (high sensitivity), may increase false positives (lower precision)
			\item To avoid false positives (high precision), may miss positives (lower sensitivity)
			\item Example: Fraud detection - high sensitivity vs. precision
			\item Choice depends on context
		\end{itemize}
		
		\vspace{0.1cm}
		\textcolor{myred}{\textbf{WHY OTHER OPTIONS ARE WRONG:}}
		\begin{itemize}
			\item \textcolor{myred}{A:} They often move in opposite directions
			\item \textcolor{myred}{C:} They are related through the threshold
			\item \textcolor{myred}{D:} Importance depends on context
		\end{itemize}
		
		\textcolor{myred}{\textbf{EXAM TRAP:}} Precision and sensitivity have a \textcolor{myblue}{trade-off}!
	\end{frame}
	
	\begin{frame}{Q24: False Positive Rate Definition - Tutorial}
		\fontsize{6.5}{7.5}\selectfont
		\textbf{QUESTION: What is the False Positive Rate (FPR)?}
		
		\vspace{0.1cm}
		\textbf{OPTIONS:}
		\begin{enumerate}
			\item FP/(TP+FP)
			\item FP/(FP+TN)
			\item FN/(TP+FN)
			\item TN/(TN+FP)
		\end{enumerate}
		
		\vspace{0.1cm}
		\textcolor{myblue}{\textbf{ANSWER: B}}
		
		\vspace{0.1cm}
		\textbf{DETAILED EXPLANATION:}
		\begin{itemize}
			\item Proportion of actual negatives that are incorrectly classified as positive
			\item $FPR = \frac{FP}{FP+TN}$
			\item Also equals 1 - Specificity
			\item Important for ROC curves and fairness analysis
		\end{itemize}
		
		\vspace{0.1cm}
		\textcolor{myred}{\textbf{WHY OTHER OPTIONS ARE WRONG:}}
		\begin{itemize}
			\item \textcolor{myred}{A:} This is 1 - precision, not FPR
			\item \textcolor{myred}{C:} This is False Negative Rate
			\item \textcolor{myred}{D:} This is specificity
		\end{itemize}
		
		\textcolor{myred}{\textbf{EXAM TRAP:}} FPR = \textcolor{myblue}{FP/(FP+TN)}!
	\end{frame}
	
	\begin{frame}{Q25: Confusion Matrix Components - Tutorial}
		\fontsize{6.5}{7.5}\selectfont
		\textbf{QUESTION: What are the four components of a confusion matrix?}
		
		\vspace{0.1cm}
		\textbf{OPTIONS:}
		\begin{enumerate}
			\item TP, TN, FP, FN
			\item Mean, Median, Mode, Standard Deviation
			\item Training, Validation, Test, Production
			\item Precision, Recall, Accuracy, F1
		\end{enumerate}
		
		\vspace{0.1cm}
		\textcolor{myblue}{\textbf{ANSWER: A}}
		
		\vspace{0.1cm}
		\textbf{DETAILED EXPLANATION:}
		\begin{itemize}
			\item \textbf{TP (True Positive):} Correctly predicted positive
			\item \textbf{TN (True Negative):} Correctly predicted negative
			\item \textbf{FP (False Positive):} Incorrectly predicted positive
			\item \textbf{FN (False Negative):} Incorrectly predicted negative
		\end{itemize}
		
		\vspace{0.1cm}
		\textcolor{myred}{\textbf{WHY OTHER OPTIONS ARE WRONG:}}
		\begin{itemize}
			\item \textcolor{myred}{B:} These are statistical measures
			\item \textcolor{myred}{C:} These are data splits
			\item \textcolor{myred}{D:} These are performance metrics derived from confusion matrix
		\end{itemize}
		
		\textcolor{myred}{\textbf{EXAM TRAP:}} Confusion matrix = \textcolor{myblue}{TP, TN, FP, FN}!
	\end{frame}
	
	\begin{frame}{Q26: Fairness Measures Summary Table - Tutorial}
		\fontsize{6.5}{7.5}\selectfont
		\textbf{QUESTION: What question does demographic parity answer?}
		
		\vspace{0.1cm}
		\textbf{OPTIONS:}
		\begin{enumerate}
			\item "Is the likelihood of a true positive prediction the same across groups?"
			\item "Is the likelihood of a positive prediction the same across groups?"
			\item "Is the likelihood of identifying actual positives as positive the same across groups?"
			\item "Are all predictions accurate?"
		\end{enumerate}
		
		\vspace{0.1cm}
		\textcolor{myblue}{\textbf{ANSWER: B}}
		
		\vspace{0.1cm}
		\textbf{DETAILED EXPLANATION:}
		\begin{itemize}
			\item Demographic parity: equal probability of positive prediction
			\item $P(\hat{Y}=1|A=0) = P(\hat{Y}=1|A=1)$
			\item Independent of whether predictions are correct
			\item Focuses on outcome distribution
		\end{itemize}
		
		\vspace{0.1cm}
		\textcolor{myred}{\textbf{WHY OTHER OPTIONS ARE WRONG:}}
		\begin{itemize}
			\item \textcolor{myred}{A:} This is predictive rate parity
			\item \textcolor{myred}{C:} This is equal opportunity
			\item \textcolor{myred}{D:} This is accuracy
		\end{itemize}
		
		\textcolor{myred}{\textbf{EXAM TRAP:}} Demographic parity = \textcolor{myblue}{equal positive prediction probability}!
	\end{frame}
	
	\begin{frame}{Q27: Fairness Measures Summary Table 2 - Tutorial}
		\fontsize{6.5}{7.5}\selectfont
		\textbf{QUESTION: What question does equal opportunity answer?}
		
		\vspace{0.1cm}
		\textbf{OPTIONS:}
		\begin{enumerate}
			\item "Is the likelihood of a positive prediction the same across groups?"
			\item "Is the likelihood of a true positive prediction the same across groups?"
			\item "Is the likelihood of identifying actual positives as positive the same across groups?"
			\item "Is the likelihood of a negative prediction the same across groups?"
		\end{enumerate}
		
		\vspace{0.1cm}
		\textcolor{myblue}{\textbf{ANSWER: C}}
		
		\vspace{0.1cm}
		\textbf{DETAILED EXPLANATION:}
		\begin{itemize}
			\item Equal opportunity: equal true positive rate
			\item $P(\hat{Y}=1|Y=1,A=0) = P(\hat{Y}=1|Y=1,A=1)$
			\item Ensures actual positives are equally likely to be caught
			\item Important when missing positives is costly
		\end{itemize}
		
		\vspace{0.1cm}
		\textcolor{myred}{\textbf{WHY OTHER OPTIONS ARE WRONG:}}
		\begin{itemize}
			\item \textcolor{myred}{A:} This is demographic parity
			\item \textcolor{myred}{B:} This is predictive rate parity
			\item \textcolor{myred}{D:} This is not a standard fairness measure
		\end{itemize}
		
		\textcolor{myred}{\textbf{EXAM TRAP:}} Equal opportunity = \textcolor{myblue}{equal true positive rate}!
	\end{frame}
	
	\begin{frame}{Q28: Base Rate Definition - Tutorial}
		\fontsize{6.5}{7.5}\selectfont
		\textbf{QUESTION: What is a base rate?}
		
		\vspace{0.1cm}
		\textbf{OPTIONS:}
		\begin{enumerate}
			\item The error rate of an algorithm
			\item The general prevalence of something within a population
			\item The accuracy of predictions
			\item The precision of predictions
		\end{enumerate}
		
		\vspace{0.1cm}
		\textcolor{myblue}{\textbf{ANSWER: B}}
		
		\vspace{0.1cm}
		\textbf{DETAILED EXPLANATION:}
		\begin{itemize}
			\item Example: Prevalence of a disease in a population
			\item Base rates can differ across groups
			\item Important for fairness analysis
			\item Different base rates create fairness trade-offs
		\end{itemize}
		
		\vspace{0.1cm}
		\textcolor{myred}{\textbf{WHY OTHER OPTIONS ARE WRONG:}}
		\begin{itemize}
			\item \textcolor{myred}{A:} Error rate is different
			\item \textcolor{myred}{C:} Accuracy is different
			\item \textcolor{myred}{D:} Precision is different
		\end{itemize}
		
		\textcolor{myred}{\textbf{EXAM TRAP:}} Base rate = \textcolor{myblue}{prevalence in a population}!
	\end{frame}
	
	\begin{frame}{Q29: Performance vs Fairness Trade-off - Tutorial}
		\fontsize{6.5}{7.5}\selectfont
		\textbf{QUESTION: Why is balancing performance and fairness challenging?}
		
		\vspace{0.1cm}
		\textbf{OPTIONS:}
		\begin{enumerate}
			\item Fairness always improves performance
			\item Optimizing for fairness can degrade performance and vice versa
			\item They are always perfectly aligned
			\item Performance is never affected by fairness
		\end{enumerate}
		
		\vspace{0.1cm}
		\textcolor{myblue}{\textbf{ANSWER: B}}
		
		\vspace{0.1cm}
		\textbf{DETAILED EXPLANATION:}
		\begin{itemize}
			\item Example: Fraud detection
			\item High sensitivity may lower precision
			\item Fairness constraints may reduce overall accuracy
			\item Need to balance based on context
		\end{itemize}
		
		\vspace{0.1cm}
		\textcolor{myred}{\textbf{WHY OTHER OPTIONS ARE WRONG:}}
		\begin{itemize}
			\item \textcolor{myred}{A:} Fairness often TRADES OFF with performance
			\item \textcolor{myred}{C:} They are often in tension
			\item \textcolor{myred}{D:} Performance IS affected by fairness constraints
		\end{itemize}
		
		\textcolor{myred}{\textbf{EXAM TRAP:}} There is often a \textcolor{myblue}{trade-off between performance and fairness}!
	\end{frame}
	
	\begin{frame}{Q30: Fairness Measures Summary - Tutorial}
		\fontsize{6.5}{7.5}\selectfont
		\textbf{QUESTION: Which fairness measure is most appropriate for medical diagnostics?}
		
		\vspace{0.1cm}
		\textbf{OPTIONS:}
		\begin{enumerate}
			\item Demographic parity
			\item Predictive rate parity
			\item Equal opportunity
			\item Accuracy parity
		\end{enumerate}
		
		\vspace{0.1cm}
		\textcolor{myblue}{\textbf{ANSWER: C}}
		
		\vspace{0.1cm}
		\textbf{DETAILED EXPLANATION:}
		\begin{itemize}
			\item Equal opportunity ensures positive cases are caught
			\item Missing a disease (false negative) can be life-threatening
			\item Same TPR across groups is critical
			\item Focuses on sensitivity (not missing positives)
		\end{itemize}
		
		\vspace{0.1cm}
		\textcolor{myred}{\textbf{WHY OTHER OPTIONS ARE WRONG:}}
		\begin{itemize}
			\item \textcolor{myred}{A:} Demographic parity ignores accuracy
			\item \textcolor{myred}{B:} Predictive rate parity focuses on positive predictions
			\item \textcolor{myred}{D:} Accuracy may be misleading with imbalanced data
		\end{itemize}
		
		\textcolor{myred}{\textbf{EXAM TRAP:}} Medical diagnostics prioritizes \textcolor{myblue}{equal opportunity} (don't miss disease)!
	\end{frame}
	
	% ============================================================
	% SECTION 3: SOURCES OF UNFAIRNESS - QUESTIONS 31-50
	% ============================================================
	
	\begin{frame}{Q31: Problem Specification Bias - Tutorial}
		\fontsize{6.5}{7.5}\selectfont
		\textbf{QUESTION: Why can problem specification introduce bias?}
		
		\vspace{0.1cm}
		\textbf{OPTIONS:}
		\begin{enumerate}
			\item It always produces perfect algorithms
			\item How a problem is specified affects algorithmic performance and fairness
			\item Problem specification never affects outcomes
			\item All problems are specified identically
		\end{enumerate}
		
		\vspace{0.1cm}
		\textcolor{myblue}{\textbf{ANSWER: B}}
		
		\vspace{0.1cm}
		\textbf{DETAILED EXPLANATION:}
		\begin{itemize}
			\item Choice of criteria affects who applies and who gets hired
			\item Translation to measurable features requires proxies
			\item Some concepts are harder to quantify than others
			\item Example: Productivity vs. teamwork
		\end{itemize}
		
		\vspace{0.1cm}
		\textcolor{myred}{\textbf{WHY OTHER OPTIONS ARE WRONG:}}
		\begin{itemize}
			\item \textcolor{myred}{A:} Problem specification can introduce bias
			\item \textcolor{myred}{C:} Problem specification significantly affects outcomes
			\item \textcolor{myred}{D:} Problems vary widely in specification
		\end{itemize}
		
		\textcolor{myred}{\textbf{EXAM TRAP:}} Problem specification = \textcolor{myblue}{critical for fairness}!
	\end{frame}
	
	\begin{frame}{Q32: Proxy Variables - Tutorial}
		\fontsize{6.5}{7.5}\selectfont
		\textbf{QUESTION: What is a proxy variable in algorithmic design?}
		
		\vspace{0.1cm}
		\textbf{OPTIONS:}
		\begin{enumerate}
			\item A variable that perfectly measures the desired concept
			\item A variable used as a stand-in for a concept that is hard to measure directly
			\item A variable that is never used in algorithms
			\item A variable with no correlation to the desired concept
		\end{enumerate}
		
		\vspace{0.1cm}
		\textcolor{myblue}{\textbf{ANSWER: B}}
		
		\vspace{0.1cm}
		\textbf{DETAILED EXPLANATION:}
		\begin{itemize}
			\item Used when direct measurement is difficult
			\item Example: Productivity measured by sales
			\item Risk: Proxies reduce complex concepts
			\item Can introduce bias if proxies are imperfect
		\end{itemize}
		
		\vspace{0.1cm}
		\textcolor{myred}{\textbf{WHY OTHER OPTIONS ARE WRONG:}}
		\begin{itemize}
			\item \textcolor{myred}{A:} Proxies are imperfect by definition
			\item \textcolor{myred}{C:} Proxies ARE used in algorithms
			\item \textcolor{myred}{D:} Proxies should be correlated with the concept
		\end{itemize}
		
		\textcolor{myred}{\textbf{EXAM TRAP:}} Proxies are \textcolor{myblue}{imperfect stand-ins} for hard-to-measure concepts!
	\end{frame}
	
	\begin{frame}{Q33: Amazon Recruiting Tool Example - Tutorial}
		\fontsize{6.5}{7.5}\selectfont
		\textbf{QUESTION: What happened with Amazon's recruiting tool?}
		
		\vspace{0.1cm}
		\textbf{OPTIONS:}
		\begin{enumerate}
			\item It worked perfectly and was deployed
			\item It systematically downgraded CVs from women applicants
			\item It eliminated all bias
			\item It was immediately successful
		\end{enumerate}
		
		\vspace{0.1cm}
		\textcolor{myblue}{\textbf{ANSWER: B}}
		
		\vspace{0.1cm}
		\textbf{DETAILED EXPLANATION:}
		\begin{itemize}
			\item Amazon attempted to create an unbiased recruiting tool
			\item Algorithm systematically downgraded women's CVs
			\item Even identical CVs were rated lower for women
			\item Amazon eventually abandoned the project
		\end{itemize}
		
		\vspace{0.1cm}
		\textcolor{myred}{\textbf{WHY OTHER OPTIONS ARE WRONG:}}
		\begin{itemize}
			\item \textcolor{myred}{A:} The tool failed due to bias
			\item \textcolor{myred}{C:} It perpetuated existing biases
			\item \textcolor{myred}{D:} The project was abandoned
		\end{itemize}
		
		\textcolor{myred}{\textbf{EXAM TRAP:}} Amazon's tool \textcolor{myblue}{systematically discriminated against women}!
	\end{frame}
	
	\begin{frame}{Q34: Fairness Through Unawareness - Tutorial}
		\fontsize{6.5}{7.5}\selectfont
		\textbf{QUESTION: What is "fairness through unawareness"?}
		
		\vspace{0.1cm}
		\textbf{OPTIONS:}
		\begin{enumerate}
			\item Using all available features
			\item Removing protected characteristics from input features
			\item Always including protected characteristics
			\item Ignoring all features
		\end{enumerate}
		
		\vspace{0.1cm}
		\textcolor{myblue}{\textbf{ANSWER: B}}
		
		\vspace{0.1cm}
		\textbf{DETAILED EXPLANATION:}
		\begin{itemize}
			\item Intuitive approach: don't use protected characteristics
			\item Assumes removing them prevents discrimination
			\item Problem: correlated variables can still carry protected information
			\item Example: Prior convictions correlated with protected group membership
		\end{itemize}
		
		\vspace{0.1cm}
		\textcolor{myred}{\textbf{WHY OTHER OPTIONS ARE WRONG:}}
		\begin{itemize}
			\item \textcolor{myred}{A:} This would increase bias risk
			\item \textcolor{myred}{C:} This would allow explicit discrimination
			\item \textcolor{myred}{D:} This would make predictions impossible
		\end{itemize}
		
		\textcolor{myred}{\textbf{EXAM TRAP:}} Fairness through unawareness \textcolor{myblue}{does NOT eliminate bias}!
	\end{frame}
	
	\begin{frame}{Q35: FTA Prediction Example - Tutorial}
		\fontsize{6.5}{7.5}\selectfont
		\textbf{QUESTION: In the FTA prediction example, what happened when protected characteristics were removed?}
		
		\vspace{0.1cm}
		\textbf{OPTIONS:}
		\begin{enumerate}
			\item Bias was completely eliminated
			\item The coefficient on prior convictions increased from 54% to 72%
			\item The algorithm became perfectly fair
			\item No changes occurred
		\end{enumerate}
		
		\vspace{0.1cm}
		\textcolor{myblue}{\textbf{ANSWER: B}}
		
		\vspace{0.1cm}
		\textbf{DETAILED EXPLANATION:}
		\begin{itemize}
			\item Protected and prior convictions were correlated
			\item Removing protected variable contaminated the prior variable
			\item Prior conviction now carried more weight
			\item Algorithm still indirectly discriminated
		\end{itemize}
		
		\vspace{0.1cm}
		\textcolor{myred}{\textbf{WHY OTHER OPTIONS ARE WRONG:}}
		\begin{itemize}
			\item \textcolor{myred}{A:} Bias was NOT eliminated
			\item \textcolor{myred}{C:} The algorithm remained biased
			\item \textcolor{myred}{D:} Significant changes occurred
		\end{itemize}
		
		\textcolor{myred}{\textbf{EXAM TRAP:}} Removing protected characteristics \textcolor{myblue}{can make bias worse} through correlated variables!
	\end{frame}
	
	\begin{frame}{Q36: Historical Bias - Tutorial}
		\fontsize{6.5}{7.5}\selectfont
		\textbf{QUESTION: What is historical bias in AI?}
		
		\vspace{0.1cm}
		\textbf{OPTIONS:}
		\begin{enumerate}
			\item Bias that occurs only in the present
			\item Bias in training data that reflects past discrimination
			\item Bias that never affects algorithms
			\item Bias that is always intentional
		\end{enumerate}
		
		\vspace{0.1cm}
		\textcolor{myblue}{\textbf{ANSWER: B}}
		
		\vspace{0.1cm}
		\textbf{DETAILED EXPLANATION:}
		\begin{itemize}
			\item Supervised algorithms learn from historical data
			\item Historical data reflects past biases and discrimination
			\item Example: Women historically less likely selected for roles
			\item Algorithm repeats and may exacerbate patterns
		\end{itemize}
		
		\vspace{0.1cm}
		\textcolor{myred}{\textbf{WHY OTHER OPTIONS ARE WRONG:}}
		\begin{itemize}
			\item \textcolor{myred}{A:} Historical bias is about the PAST
			\item \textcolor{myred}{C:} Historical bias DOES affect algorithms
			\item \textcolor{myred}{D:} Historical bias is often unintentional
		\end{itemize}
		
		\textcolor{myred}{\textbf{EXAM TRAP:}} Historical bias = \textcolor{myblue}{past discrimination reflected in training data}!
	\end{frame}
	
	\begin{frame}{Q37: Feedback Loops - Tutorial}
		\fontsize{6.5}{7.5}\selectfont
		\textbf{QUESTION: What is a feedback loop in AI systems?}
		
		\vspace{0.1cm}
		\textbf{OPTIONS:}
		\begin{enumerate}
			\item A system that corrects its own errors
			\item When algorithmic predictions influence future data, perpetuating bias
			\item A system that never changes
			\item A system that eliminates all bias
		\end{enumerate}
		
		\vspace{0.1cm}
		\textcolor{myblue}{\textbf{ANSWER: B}}
		
		\vspace{0.1cm}
		\textbf{DETAILED EXPLANATION:}
		\begin{itemize}
			\item Algorithm's outputs influence future data
			\item Biased predictions create more biased data
			\item Amplifies existing discrimination
			\item Example: Hiring algorithm learns from its own biased outcomes
		\end{itemize}
		
		\vspace{0.1cm}
		\textcolor{myred}{\textbf{WHY OTHER OPTIONS ARE WRONG:}}
		\begin{itemize}
			\item \textcolor{myred}{A:} Feedback loops often AMPLIFY errors
			\item \textcolor{myred}{C:} Feedback loops cause change (often negative)
			\item \textcolor{myred}{D:} Feedback loops can perpetuate bias
		\end{itemize}
		
		\textcolor{myred}{\textbf{EXAM TRAP:}} Feedback loops \textcolor{myblue}{amplify and perpetuate bias}!
	\end{frame}
	
	\begin{frame}{Q38: Representation Bias - Tutorial}
		\fontsize{6.5}{7.5}\selectfont
		\textbf{QUESTION: What is representation bias?}
		
		\vspace{0.1cm}
		\textbf{OPTIONS:}
		\begin{enumerate}
			\item When data is perfectly representative
			\item When the collected sample is not representative of the target population
			\item When all groups are equally represented
			\item When data collection is unbiased
		\end{enumerate}
		
		\vspace{0.1cm}
		\textcolor{myblue}{\textbf{ANSWER: B}}
		
		\vspace{0.1cm}
		\textbf{DETAILED EXPLANATION:}
		\begin{itemize}
			\item Sample does not reflect target population
			\item Results in problematic model specifications
			\item Example: Facial recognition trained mostly on white faces
			\item Algorithm performs better on dominant groups
		\end{itemize}
		
		\vspace{0.1cm}
		\textcolor{myred}{\textbf{WHY OTHER OPTIONS ARE WRONG:}}
		\begin{itemize}
			\item \textcolor{myred}{A:} Representation bias means NOT representative
			\item \textcolor{myred}{C:} This would be balanced representation
			\item \textcolor{myred}{D:} Representation bias IS a bias
		\end{itemize}
		
		\textcolor{myred}{\textbf{EXAM TRAP:}} Representation bias = \textcolor{myblue}{non-representative sample}!
	\end{frame}
	
	\begin{frame}{Q39: Measurement Bias - Tutorial}
		\fontsize{6.5}{7.5}\selectfont
		\textbf{QUESTION: What is measurement bias?}
		
		\vspace{0.1cm}
		\textbf{OPTIONS:}
		\begin{enumerate}
			\item When measurement methods are consistent across the sample
			\item When measurement methods are not consistent across the sample
			\item When all measurements are perfectly accurate
			\item When measurements are always reliable
		\end{enumerate}
		
		\vspace{0.1cm}
		\textcolor{myblue}{\textbf{ANSWER: B}}
		
		\vspace{0.1cm}
		\textbf{DETAILED EXPLANATION:}
		\begin{itemize}
			\item Inconsistent measurement methods
			\item Can favor certain groups
			\item Example: Doctor spends more time with wealthy patients
			\item Data on early-stage disease discovery may be biased
		\end{itemize}
		
		\vspace{0.1cm}
		\textcolor{myred}{\textbf{WHY OTHER OPTIONS ARE WRONG:}}
		\begin{itemize}
			\item \textcolor{myred}{A:} Measurement bias means INCONSISTENT methods
			\item \textcolor{myred}{C:} Measurement bias means NOT perfectly accurate
			\item \textcolor{myred}{D:} Measurement bias means UNRELIABLE across groups
		\end{itemize}
		
		\textcolor{myred}{\textbf{EXAM TRAP:}} Measurement bias = \textcolor{myblue}{inconsistent measurement across groups}!
	\end{frame}
	
	\begin{frame}{Q40: Data Composition Bias - Tutorial}
		\fontsize{6.5}{7.5}\selectfont
		\textbf{QUESTION: What is data composition bias?}
		
		\vspace{0.1cm}
		\textbf{OPTIONS:}
		\begin{enumerate}
			\item When datasets are perfectly balanced
			\item When minority subgroups are underrepresented, causing algorithms to underperform
			\item When all groups are equally represented
			\item When data is always complete
		\end{enumerate}
		
		\vspace{0.1cm}
		\textcolor{myblue}{\textbf{ANSWER: B}}
		
		\vspace{0.1cm}
		\textbf{DETAILED EXPLANATION:}
		\begin{itemize}
			\item Statistical minorities have fewer data points
			\item Algorithm underperforms for underrepresented groups
			\item Example: 3% pregnant women in medical dataset
			\item Algorithm may significantly underperform for this group
		\end{itemize}
		
		\vspace{0.1cm}
		\textcolor{myred}{\textbf{WHY OTHER OPTIONS ARE WRONG:}}
		\begin{itemize}
			\item \textcolor{myred}{A:} Composition bias means IMBALANCED data
			\item \textcolor{myred}{C:} This would be balanced data
			\item \textcolor{myred}{D:} Composition bias means INCOMPLETE representation
		\end{itemize}
		
		\textcolor{myred}{\textbf{EXAM TRAP:}} Data composition bias = \textcolor{myblue}{underrepresented minority subgroups}!
	\end{frame}
	
	\begin{frame}{Q41: Balancing Datasets - Tutorial}
		\fontsize{6.5}{7.5}\selectfont
		\textbf{QUESTION: What does balancing a dataset mean?}
		
		\vspace{0.1cm}
		\textbf{OPTIONS:}
		\begin{enumerate}
			\item Using proportional representation
			\item Ensuring relevant minority subgroups occupy a similar proportion as majority groups
			\item Removing all minority data
			\item Always using the original distribution
		\end{enumerate}
		
		\vspace{0.1cm}
		\textcolor{myblue}{\textbf{ANSWER: B}}
		
		\vspace{0.1cm}
		\textbf{DETAILED EXPLANATION:}
		\begin{itemize}
			\item Minority groups represented more than proportionally
			\item Ensures algorithm recognizes minority groups
			\item Different from proportional representation
			\item Example: 50% male, 50% female even if 80/20 historically
		\end{itemize}
		
		\vspace{0.1cm}
		\textcolor{myred}{\textbf{WHY OTHER OPTIONS ARE WRONG:}}
		\begin{itemize}
			\item \textcolor{myred}{A:} Proportional representation is different
			\item \textcolor{myred}{C:} Balancing increases minority representation
			\item \textcolor{myred}{D:} Balancing changes the original distribution
		\end{itemize}
		
		\textcolor{myred}{\textbf{EXAM TRAP:}} Balancing = \textcolor{myblue}{over-representing minorities} for algorithm performance!
	\end{frame}
	
	\begin{frame}{Q42: Bootstrapping Methods - Tutorial}
		\fontsize{6.5}{7.5}\selectfont
		\textbf{QUESTION: What are methods for balancing datasets?}
		
		\vspace{0.1cm}
		\textbf{OPTIONS:}
		\begin{enumerate}
			\item Over-sampling minority groups
			\item Under-sampling majority groups
			\item Both A and B
			\item Only removing outliers
		\end{enumerate}
		
		\vspace{0.1cm}
		\textcolor{myblue}{\textbf{ANSWER: C}}
		
		\vspace{0.1cm}
		\textbf{DETAILED EXPLANATION:}
		\begin{itemize}
			\item \textbf{Over-sampling:} Add more data from minority group
			\item \textbf{Under-sampling:} Remove data from majority group
			\item \textbf{Combination:} Use both methods
			\item Can also create synthetic data points
		\end{itemize}
		
		\vspace{0.1cm}
		\textcolor{myred}{\textbf{WHY OTHER OPTIONS ARE WRONG:}}
		\begin{itemize}
			\item \textcolor{myred}{A:} Correct but incomplete
			\item \textcolor{myred}{B:} Correct but incomplete
			\item \textcolor{myred}{D:} Outlier removal is different
		\end{itemize}
		
		\textcolor{myred}{\textbf{EXAM TRAP:}} Bootstrapping methods include \textcolor{myblue}{over-sampling AND under-sampling}!
	\end{frame}
	
	\begin{frame}{Q43: Intersectional Bias - Tutorial}
		\fontsize{6.5}{7.5}\selectfont
		\textbf{QUESTION: What is intersectional bias?}
		
		\vspace{0.1cm}
		\textbf{OPTIONS:}
		\begin{enumerate}
			\item Bias that only affects one group
			\item Bias affecting groups at the intersection of multiple underrepresented subgroups
			\item Bias that never occurs
			\item Bias that is always intentional
		\end{enumerate}
		
		\vspace{0.1cm}
		\textcolor{myblue}{\textbf{ANSWER: B}}
		
		\vspace{0.1cm}
		\textbf{DETAILED EXPLANATION:}
		\begin{itemize}
			\item Performance can be particularly low for intersectional groups
			\item Example: Women of color in facial recognition
			\item Multiple dimensions of underrepresentation
			\item Harder to detect and correct
		\end{itemize}
		
		\vspace{0.1cm}
		\textcolor{myred}{\textbf{WHY OTHER OPTIONS ARE WRONG:}}
		\begin{itemize}
			\item \textcolor{myred}{A:} Intersectional bias affects MULTIPLE groups
			\item \textcolor{myred}{C:} Intersectional bias DOES occur
			\item \textcolor{myred}{D:} Intersectional bias is often unintentional
		\end{itemize}
		
		\textcolor{myred}{\textbf{EXAM TRAP:}} Intersectional bias = \textcolor{myblue}{multiple underrepresented dimensions combined}!
	\end{frame}
	
	\begin{frame}{Q44: Word Embeddings Bias - Tutorial}
		\fontsize{6.5}{7.5}\selectfont
		\textbf{QUESTION: How can word embeddings introduce bias?}
		
		\vspace{0.1cm}
		\textbf{OPTIONS:}
		\begin{enumerate}
			\item They always produce unbiased results
			\item They can contain stereotypes (e.g., man:computer programmer as woman:homemaker)
			\item They never contain stereotypes
			\item They are always perfectly neutral
		\end{enumerate}
		
		\vspace{0.1cm}
		\textcolor{myblue}{\textbf{ANSWER: B}}
		
		\vspace{0.1cm}
		\textbf{DETAILED EXPLANATION:}
		\begin{itemize}
			\item Word embeddings capture meaning from text
			\item Text contains societal stereotypes
			\item Example: "man is to computer programmer as woman is to homemaker"
			\item NLP algorithms can incorporate these stereotypes
		\end{itemize}
		
		\vspace{0.1cm}
		\textcolor{myred}{\textbf{WHY OTHER OPTIONS ARE WRONG:}}
		\begin{itemize}
			\item \textcolor{myred}{A:} Word embeddings CAN contain bias
			\item \textcolor{myred}{C:} Word embeddings DO contain stereotypes
			\item \textcolor{myred}{D:} Word embeddings reflect societal biases
		\end{itemize}
		
		\textcolor{myred}{\textbf{EXAM TRAP:}} Word embeddings can contain \textcolor{myblue}{harmful stereotypes}!
	\end{frame}
	
	\begin{frame}{Q45: Objective Function Bias - Tutorial}
		\fontsize{6.5}{7.5}\selectfont
		\textbf{QUESTION: How can objective functions introduce bias?}
		
		\vspace{0.1cm}
		\textbf{OPTIONS:}
		\begin{enumerate}
			\item They never affect fairness
			\item What the algorithm optimizes for makes an ethically significant difference
			\item All objective functions are equally fair
			\item Objective functions always eliminate bias
		\end{enumerate}
		
		\vspace{0.1cm}
		\textcolor{myblue}{\textbf{ANSWER: B}}
		
		\vspace{0.1cm}
		\textbf{DETAILED EXPLANATION:}
		\begin{itemize}
			\item Focus on performance vs. error distribution
			\item Cancer detection example: err on side of caution
			\item Brain cancer: invasive treatment considerations
			\item Choice depends on context
		\end{itemize}
		
		\vspace{0.1cm}
		\textcolor{myred}{\textbf{WHY OTHER OPTIONS ARE WRONG:}}
		\begin{itemize}
			\item \textcolor{myred}{A:} Objective functions significantly affect fairness
			\item \textcolor{myred}{C:} Different objectives have different fairness implications
			\item \textcolor{myred}{D:} Objective functions CAN introduce bias
		\end{itemize}
		
		\textcolor{myred}{\textbf{EXAM TRAP:}} Objective function choice = \textcolor{myblue}{ethical significance}!
	\end{frame}
	
	\begin{frame}{Q46: Cancer Detection Example - Tutorial}
		\fontsize{6.5}{7.5}\selectfont
		\textbf{QUESTION: In the cancer detection example, why might different algorithms require different approaches?}
		
		\vspace{0.1cm}
		\textbf{OPTIONS:}
		\begin{enumerate}
			\item All cancers are the same
			\item Skin cancer and brain cancer have different treatment risks and consequences
			\item Treatment is never invasive
			\item All patients are identical
		\end{enumerate}
		
		\vspace{0.1cm}
		\textcolor{myblue}{\textbf{ANSWER: B}}
		
		\vspace{0.1cm}
		\textbf{DETAILED EXPLANATION:}
		\begin{itemize}
			\item Skin cancer: biopsy, relatively low risk
			\item Brain cancer: major surgery, high risk
			\item Different trade-offs for false positives/negatives
			\item Context matters for objective function
		\end{itemize}
		
		\vspace{0.1cm}
		\textcolor{myred}{\textbf{WHY OTHER OPTIONS ARE WRONG:}}
		\begin{itemize}
			\item \textcolor{myred}{A:} Different cancers have different risks
			\item \textcolor{myred}{C:} Some treatments ARE highly invasive
			\item \textcolor{myred}{D:} Patients have different conditions
		\end{itemize}
		
		\textcolor{myred}{\textbf{EXAM TRAP:}} Context determines \textcolor{myblue}{appropriate trade-offs}!
	\end{frame}
	
	\begin{frame}{Q47: Benchmark Dataset Bias - Tutorial}
		\fontsize{6.5}{7.5}\selectfont
		\textbf{QUESTION: How can benchmark datasets introduce bias?}
		
		\vspace{0.1cm}
		\textbf{OPTIONS:}
		\begin{enumerate}
			\item They are always representative
			\item They may not be representative of the intended use population
			\item They always eliminate bias
			\item They never affect model performance
		\end{enumerate}
		
		\vspace{0.1cm}
		\textcolor{myblue}{\textbf{ANSWER: B}}
		
		\vspace{0.1cm}
		\textbf{DETAILED EXPLANATION:}
		\begin{itemize}
			\item Public benchmarks may not reflect real-world population
			\item Good performance on benchmark doesn't guarantee real-world performance
			\item Can mask biases that appear in deployment
			\item Need domain-specific testing
		\end{itemize}
		
		\vspace{0.1cm}
		\textcolor{myred}{\textbf{WHY OTHER OPTIONS ARE WRONG:}}
		\begin{itemize}
			\item \textcolor{myred}{A:} Benchmarks may NOT be representative
			\item \textcolor{myred}{C:} Benchmarks can CONTAIN bias
			\item \textcolor{myred}{D:} Benchmarks significantly affect perceived performance
		\end{itemize}
		
		\textcolor{myred}{\textbf{EXAM TRAP:}} Benchmark datasets may \textcolor{myblue}{not reflect real-world populations}!
	\end{frame}
	
	\begin{frame}{Q48: Deployment Context Bias - Tutorial}
		\fontsize{6.5}{7.5}\selectfont
		\textbf{QUESTION: How can deployment context introduce bias?}
		
		\vspace{0.1cm}
		\textbf{OPTIONS:}
		\begin{enumerate}
			\item Algorithms always perform equally in all contexts
			\item Algorithms used outside their designed context can lead to performance loss
			\item Context never matters
			\item All contexts are identical
		\end{enumerate}
		
		\vspace{0.1cm}
		\textcolor{myblue}{\textbf{ANSWER: B}}
		
		\vspace{0.1cm}
		\textbf{DETAILED EXPLANATION:}
		\begin{itemize}
			\item Algorithm for one context may not work in another
			\item Example: Recidivism prediction for sentence length
			\item Historical data may be outdated
			\item Values of users may differ from developers
		\end{itemize}
		
		\vspace{0.1cm}
		\textcolor{myred}{\textbf{WHY OTHER OPTIONS ARE WRONG:}}
		\begin{itemize}
			\item \textcolor{myred}{A:} Performance varies by context
			\item \textcolor{myred}{C:} Context significantly matters
			\item \textcolor{myred}{D:} Contexts vary widely
		\end{itemize}
		
		\textcolor{myred}{\textbf{EXAM TRAP:}} Deployment context = \textcolor{myblue}{critical for performance and fairness}!
	\end{frame}
	
	\begin{frame}{Q49: Human Oversight Importance - Tutorial}
		\fontsize{6.5}{7.5}\selectfont
		\textbf{QUESTION: Why is human oversight important for AI systems?}
		
		\vspace{0.1cm}
		\textbf{OPTIONS:}
		\begin{enumerate}
			\item Humans are never needed
			\item Ensuring safe and fair algorithms requires human oversight throughout
			\item AI systems are always perfectly safe
			\item Oversight is optional
		\end{enumerate}
		
		\vspace{0.1cm}
		\textcolor{myblue}{\textbf{ANSWER: B}}
		
		\vspace{0.1cm}
		\textbf{DETAILED EXPLANATION:}
		\begin{itemize}
			\item Data collection process needs oversight
			\item Model development needs review
			\item Deployment needs monitoring
			\item Regular auditing becomes important
			\item Consequential domains require accountability
		\end{itemize}
		
		\vspace{0.1cm}
		\textcolor{myred}{\textbf{WHY OTHER OPTIONS ARE WRONG:}}
		\begin{itemize}
			\item \textcolor{myred}{A:} Human oversight is ESSENTIAL
			\item \textcolor{myred}{C:} AI systems are NOT always perfectly safe
			\item \textcolor{myred}{D:} Oversight is CRITICAL, not optional
		\end{itemize}
		
		\textcolor{myred}{\textbf{EXAM TRAP:}} Human oversight is \textcolor{myblue}{essential throughout the AI lifecycle}!
	\end{frame}
	
	\begin{frame}{Q50: Sources of Unfairness Summary - Tutorial}
		\fontsize{6.5}{7.5}\selectfont
		\textbf{QUESTION: Which is NOT a source of algorithmic unfairness?}
		
		\vspace{0.1cm}
		\textbf{OPTIONS:}
		\begin{enumerate}
			\item Problem specification and feature selection
			\item Data collection and composition
			\item Model development
			\item Always using perfect data
		\end{enumerate}
		
		\vspace{0.1cm}
		\textcolor{myblue}{\textbf{ANSWER: D}}
		
		\vspace{0.1cm}
		\textbf{DETAILED EXPLANATION:}
		\begin{itemize}
			\item Problem specification: choice of criteria and proxies
			\item Data collection: representation and measurement bias
			\item Data composition: imbalanced datasets
			\item Model development: embeddings, objective functions, testing
			\item Deployment: context, outdated data, user values\end{itemize}
		
		\vspace{0.1cm}
		\textcolor{myred}{\textbf{WHY OTHER OPTIONS ARE WRONG:}}
		\begin{itemize}
			\item \textcolor{myred}{A:} This IS a source of unfairness
			\item \textcolor{myred}{B:} This IS a source of unfairness
			\item \textcolor{myred}{C:} This IS a source of unfairness
		\end{itemize}
		
		\textcolor{myred}{\textbf{EXAM TRAP:}} Perfect data \textcolor{myblue}{does NOT exist} - it's a source of unfairness!
	\end{frame}
	
	% ============================================================
	% SECTION 4: EXPLAINABILITY, INTERPRETABILITY, TRANSPARENCY - QUESTIONS 51-65
	% ============================================================
	
	\begin{frame}{Q51: Explainability Definition - Tutorial}
		\fontsize{6.5}{7.5}\selectfont
		\textbf{QUESTION: What is explainability in AI?}
		
		\vspace{0.1cm}
		\textbf{OPTIONS:}
		\begin{enumerate}
			\item The ability to make predictions
			\item The capacity to explain in understandable terms how an AI system makes decisions
			\item The speed of the algorithm
			\item The accuracy of the algorithm
		\end{enumerate}
		
		\vspace{0.1cm}
		\textcolor{myblue}{\textbf{ANSWER: B}}
		
		\vspace{0.1cm}
		\textbf{DETAILED EXPLANATION:}
		\begin{itemize}
			\item Ex-post explainability (after the fact)
			\item Explaining decisions in understandable terms
			\item Essential for accountability and trust
			\item Particularly important in high-stakes domains
		\end{itemize}
		
		\vspace{0.1cm}
		\textcolor{myred}{\textbf{WHY OTHER OPTIONS ARE WRONG:}}
		\begin{itemize}
			\item \textcolor{myred}{A:} This is prediction capability
			\item \textcolor{myred}{C:} Speed is different
			\item \textcolor{myred}{D:} Accuracy is different
		\end{itemize}
		
		\textcolor{myred}{\textbf{EXAM TRAP:}} Explainability = \textcolor{myblue}{explaining decisions after the fact}!
	\end{frame}
	
	\begin{frame}{Q52: Interpretability Definition - Tutorial}
		\fontsize{6.5}{7.5}\selectfont
		\textbf{QUESTION: What is interpretability in AI?}
		
		\vspace{0.1cm}
		\textbf{OPTIONS:}
		\begin{enumerate}
			\item Explaining decisions after the fact
			\item The degree to which a human can comprehend and predict the model's behavior inherently
			\item The speed of the algorithm
			\item The accuracy of the algorithm
		\end{enumerate}
		
		\vspace{0.1cm}
		\textcolor{myblue}{\textbf{ANSWER: B}}
		
		\vspace{0.1cm}
		\textbf{DETAILED EXPLANATION:}
		\begin{itemize}
			\item Built-in mechanisms to understand inputs → outputs
			\item Inherently interpretable models
			\item Example: Decision trees (can follow the logic)
			\item Before the fact (ex-ante), not after
		\end{itemize}
		
		\vspace{0.1cm}
		\textcolor{myred}{\textbf{WHY OTHER OPTIONS ARE WRONG:}}
		\begin{itemize}
			\item \textcolor{myred}{A:} This is explainability
			\item \textcolor{myred}{C:} Speed is different
			\item \textcolor{myred}{D:} Accuracy is different
		\end{itemize}
		
		\textcolor{myred}{\textbf{EXAM TRAP:}} Interpretability = \textcolor{myblue}{inherently understandable models}!
	\end{frame}
	
	\begin{frame}{Q53: Transparency Definition - Tutorial}
		\fontsize{6.5}{7.5}\selectfont
		\textbf{QUESTION: What is transparency in AI?}
		
		\vspace{0.1cm}
		\textbf{OPTIONS:}
		\begin{enumerate}
			\item Making predictions quickly
			\item The openness and accessibility of an AI system's design, algorithms, data, and decision-making processes
			\item The accuracy of the algorithm
			\item The speed of the algorithm
		\end{enumerate}
		
		\vspace{0.1cm}
		\textcolor{myblue}{\textbf{ANSWER: B}}
		
		\vspace{0.1cm}
		\textbf{DETAILED EXPLANATION:}
		\begin{itemize}
			\item Openness about design and algorithms
			\item Accessibility of data and processes
			\item Builds trust and enables oversight
			\item Facilitates collaboration
		\end{itemize}
		
		\vspace{0.1cm}
		\textcolor{myred}{\textbf{WHY OTHER OPTIONS ARE WRONG:}}
		\begin{itemize}
			\item \textcolor{myred}{A:} Speed is different
			\item \textcolor{myred}{C:} Accuracy is different
			\item \textcolor{myred}{D:} Speed is different
		\end{itemize}
		
		\textcolor{myred}{\textbf{EXAM TRAP:}} Transparency = \textcolor{myblue}{openness and accessibility} of AI systems!
	\end{frame}
	
	\begin{frame}{Q54: Black-Box Problem - Tutorial}
		\fontsize{6.5}{7.5}\selectfont
		\textbf{QUESTION: What is the black-box problem in AI?}
		
		\vspace{0.1cm}
		\textbf{OPTIONS:}
		\begin{enumerate}
			\item Models that are always transparent
			\item Models that are too complex to be interpreted by humans
			\item Models that are always simple
			\item Models that are always accurate
		\end{enumerate}
		
		\vspace{0.1cm}
		\textcolor{myblue}{\textbf{ANSWER: B}}
		
		\vspace{0.1cm}
		\textbf{DETAILED EXPLANATION:}
		\begin{itemize}
			\item Neural networks are often black boxes
			\item Cannot understand intrinsic architecture
			\item Leads to explainability issues
			\item Particularly challenging for high-stakes decisions
		\end{itemize}
		
		\vspace{0.1cm}
		\textcolor{myred}{\textbf{WHY OTHER OPTIONS ARE WRONG:}}
		\begin{itemize}
			\item \textcolor{myred}{A:} Black boxes are NOT transparent
			\item \textcolor{myred}{C:} Black boxes are COMPLEX
			\item \textcolor{myred}{D:} Black boxes may NOT be accurate
		\end{itemize}
		
		\textcolor{myred}{\textbf{EXAM TRAP:}} Black-box problem = \textcolor{myblue}{models too complex to interpret}!
	\end{frame}
	
	\begin{frame}{Q55: Feature Importance Scores - Tutorial}
		\fontsize{6.5}{7.5}\selectfont
		\textbf{QUESTION: What are feature importance scores used for?}
		
		\vspace{0.1cm}
		\textbf{OPTIONS:}
		\begin{enumerate}
			\item To make predictions faster
			\item To rank the importance of input features for outcomes
			\item To reduce the number of data points
			\item To increase model complexity
		\end{enumerate}
		
		\vspace{0.1cm}
		\textcolor{myblue}{\textbf{ANSWER: B}}
		
		\vspace{0.1cm}
		\textbf{DETAILED EXPLANATION:}
		\begin{itemize}
			\item Used in Explainable AI (XAI)
			\item Reveals which factors have most influence
			\item Example: Credit scoring - income vs. credit history
			\item Helps identify potential biases
		\end{itemize}
		
		\vspace{0.1cm}
		\textcolor{myred}{\textbf{WHY OTHER OPTIONS ARE WRONG:}}
		\begin{itemize}
			\item \textcolor{myred}{A:} Feature importance is about interpretation
			\item \textcolor{myred}{C:} It doesn't reduce data points
			\item \textcolor{myred}{D:} It explains, not increases complexity
		\end{itemize}
		
		\textcolor{myred}{\textbf{EXAM TRAP:}} Feature importance = \textcolor{myblue}{ranks feature influence}!
	\end{frame}
	
	\begin{frame}{Q56: Shapley Values - Tutorial}
		\fontsize{6.5}{7.5}\selectfont
		\textbf{QUESTION: What do Shapley values calculate?}
		
		\vspace{0.1cm}
		\textbf{OPTIONS:}
		\begin{enumerate}
			\item The accuracy of a model
			\item The contribution of each input feature to the model's output
			\item The speed of the model
			\item The size of the model
		\end{enumerate}
		
		\vspace{0.1cm}
		\textcolor{myblue}{\textbf{ANSWER: B}}
		
		\vspace{0.1cm}
		\textbf{DETAILED EXPLANATION:}
		\begin{itemize}
			\item Calculate marginal effect of including each feature
			\item Across all possible feature combinations
			\item Fair attribution of contributions
			\item Based on game theory (Shapley values)
		\end{itemize}
		
		\vspace{0.1cm}
		\textcolor{myred}{\textbf{WHY OTHER OPTIONS ARE WRONG:}}
		\begin{itemize}
			\item \textcolor{myred}{A:} Accuracy is different
			\item \textcolor{myred}{C:} Speed is different
			\item \textcolor{myred}{D:} Size is different
		\end{itemize}
		
		\textcolor{myred}{\textbf{EXAM TRAP:}} Shapley values = \textcolor{myblue}{feature contribution calculation}!
	\end{frame}
	
	\begin{frame}{Q57: LUCID - Tutorial}
		\fontsize{6.5}{7.5}\selectfont
		\textbf{QUESTION: What does LUCID do in XAI?}
		
		\vspace{0.1cm}
		\textbf{OPTIONS:}
		\begin{enumerate}
			\item Makes predictions faster
			\item Works backward from algorithmic output to reveal internal logic
			\item Increases model complexity
			\item Removes features
		\end{enumerate}
		
		\vspace{0.1cm}
		\textcolor{myblue}{\textbf{ANSWER: B}}
		
		\vspace{0.1cm}
		\textbf{DETAILED EXPLANATION:}
		\begin{itemize}
			\item LUCID = Locating Unfairness through Canonical Inverse Design
			\item Works backward from output to input
			\item Creates distribution over input space
			\item Reveals model's internal logic
		\end{itemize}
		
		\vspace{0.1cm}
		\textcolor{myred}{\textbf{WHY OTHER OPTIONS ARE WRONG:}}
		\begin{itemize}
			\item \textcolor{myred}{A:} Speed is not LUCID's purpose
			\item \textcolor{myred}{C:} LUCID helps understand complexity
			\item \textcolor{myred}{D:} LUCID doesn't remove features
		\end{itemize}
		
		\textcolor{myred}{\textbf{EXAM TRAP:}} LUCID works \textcolor{myblue}{backward from output to reveal logic}!
	\end{frame}
	
	\begin{frame}{Q58: Surrogate Models - Tutorial}
		\fontsize{6.5}{7.5}\selectfont
		\textbf{QUESTION: What are surrogate models?}
		
		\vspace{0.1cm}
		\textbf{OPTIONS:}
		\begin{enumerate}
			\item More complex models that replace simpler ones
			\item Simpler, more interpretable models that approximate complex AI systems
			\item Models that are always less accurate
			\item Models that cannot be interpreted
		\end{enumerate}
		
		\vspace{0.1cm}
		\textcolor{myblue}{\textbf{ANSWER: B}}
		
		\vspace{0.1cm}
		\textbf{DETAILED EXPLANATION:}
		\begin{itemize}
			\item Example: LIME uses surrogate models
			\item Approximate complex model's predictions locally
			\item Use interpretable models (linear regression, decision trees)
			\item Enhance interpretability
		\end{itemize}
		
		\vspace{0.1cm}
		\textcolor{myred}{\textbf{WHY OTHER OPTIONS ARE WRONG:}}
		\begin{itemize}
			\item \textcolor{myred}{A:} Surrogate models are SIMPLER
			\item \textcolor{myred}{C:} Surrogate models can be quite accurate locally
			\item \textcolor{myred}{D:} Surrogate models are INTERPRETABLE
		\end{itemize}
		
		\textcolor{myred}{\textbf{EXAM TRAP:}} Surrogate models are \textcolor{myblue}{simpler, interpretable approximations}!
	\end{frame}
	
	\begin{frame}{Q59: LIME Definition - Tutorial}
		\fontsize{6.5}{7.5}\selectfont
		\textbf{QUESTION: What is LIME in XAI?}
		
		\vspace{0.1cm}
		\textbf{OPTIONS:}
		\begin{enumerate}
			\item Local Interpretable Model-agnostic Explanations
			\item A method to make models faster
			\item A technique to increase accuracy
			\item A method to remove bias
		\end{enumerate}
		
		\vspace{0.1cm}
		\textcolor{myblue}{\textbf{ANSWER: A}}
		
		\vspace{0.1cm}
		\textbf{DETAILED EXPLANATION:}
		\begin{itemize}
			\item LIME = Local Interpretable Model-agnostic Explanations
			\item Approximates complex model locally
			\item Uses interpretable model (linear regression, decision tree)
			\item Model-agnostic (works with any model)
		\end{itemize}
		
		\vspace{0.1cm}
		\textcolor{myred}{\textbf{WHY OTHER OPTIONS ARE WRONG:}}
		\begin{itemize}
			\item \textcolor{myred}{B:} LIME is about explanation, not speed
			\item \textcolor{myred}{C:} LIME is about interpretability, not accuracy
			\item \textcolor{myred}{D:} LIME helps identify bias, doesn't remove it directly
		\end{itemize}
		
		\textcolor{myred}{\textbf{EXAM TRAP:}} LIME = \textcolor{myblue}{Local Interpretable Model-agnostic Explanations}!
	\end{frame}
	
	\begin{frame}{Q60: XAI Challenges - Tutorial}
		\fontsize{6.5}{7.5}\selectfont
		\textbf{QUESTION: What is a challenge with XAI techniques?}
		
		\vspace{0.1cm}
		\textbf{OPTIONS:}
		\begin{enumerate}
			\item They always produce perfect explanations
			\item Oversimplification can distort understanding
			\item They are never computationally expensive
			\item They always eliminate bias
		\end{enumerate}
		
		\vspace{0.1cm}
		\textcolor{myblue}{\textbf{ANSWER: B}}
		
		\vspace{0.1cm}
		\textbf{DETAILED EXPLANATION:}
		\begin{itemize}
			\item Complex decisions oversimplified
			\item Can lead to misinterpretation
			\item Computationally intensive
			\item Ex-post techniques (after training)
			\item Trade-off: complex models vs. interpretability
		\end{itemize}
		
		\vspace{0.1cm}
		\textcolor{myred}{\textbf{WHY OTHER OPTIONS ARE WRONG:}}
		\begin{itemize}
			\item \textcolor{myred}{A:} Explanations can be imperfect
			\item \textcolor{myred}{C:} XAI CAN be computationally expensive
			\item \textcolor{myred}{D:} XAI helps identify, not eliminate bias
		\end{itemize}
		
		\textcolor{myred}{\textbf{EXAM TRAP:}} XAI can \textcolor{myblue}{oversimplify and distort understanding}!
	\end{frame}
	
	\begin{frame}{Q61: Secrecy - Tutorial}
		\fontsize{6.5}{7.5}\selectfont
		\textbf{QUESTION: What is secrecy as a form of opaqueness?}
		
		\vspace{0.1cm}
		\textbf{OPTIONS:}
		\begin{enumerate}
			\item When algorithms are fully transparent
			\item When individuals are not aware of being subject to algorithmic decision making
			\item When algorithms are always explained
			\item When all data is public
		\end{enumerate}
		
		\vspace{0.1cm}
		\textcolor{myblue}{\textbf{ANSWER: B}}
		
		\vspace{0.1cm}
		\textbf{DETAILED EXPLANATION:}
		\begin{itemize}
			\item Example: Algorithm extracts data without knowledge
			\item Used for targeted advertising
			\item Individuals unaware of algorithmic profiling
			\item Violates transparency and consent
		\end{itemize}
		
		\vspace{0.1cm}
		\textcolor{myred}{\textbf{WHY OTHER OPTIONS ARE WRONG:}}
		\begin{itemize}
			\item \textcolor{myred}{A:} Secrecy is the OPPOSITE of transparent
			\item \textcolor{myred}{C:} Secrecy means NOT explained
			\item \textcolor{myred}{D:} Secrecy means NOT public
		\end{itemize}
		
		\textcolor{myred}{\textbf{EXAM TRAP:}} Secrecy = \textcolor{myblue}{unawareness of algorithmic decision making}!
	\end{frame}
	
	\begin{frame}{Q62: Confidentiality - Tutorial}
		\fontsize{6.5}{7.5}\selectfont
		\textbf{QUESTION: What is confidentiality as a form of opaqueness?}
		
		\vspace{0.1cm}
		\textbf{OPTIONS:}
		\begin{enumerate}
			\item When algorithms are completely transparent
			\item When the algorithmic decision-making process is opaque to affected parties
			\item When all data is public
			\item When algorithms are always explained
		\end{enumerate}
		
		\vspace{0.1cm}
		\textcolor{myblue}{\textbf{ANSWER: B}}
		
		\vspace{0.1cm}
		\textbf{DETAILED EXPLANATION:}
		\begin{itemize}
			\item Individuals know they're subject to AI
			\item But don't know how decisions are made
			\item Main concern in AI discourse
			\item Similar to non-algorithmic secret decision-making
		\end{itemize}
		
		\vspace{0.1cm}
		\textcolor{myred}{\textbf{WHY OTHER OPTIONS ARE WRONG:}}
		\begin{itemize}
			\item \textcolor{myred}{A:} Confidentiality means NOT transparent
			\item \textcolor{myred}{C:} Confidentiality means NOT public
			\item \textcolor{myred}{D:} Confidentiality means NOT explained
		\end{itemize}
		
		\textcolor{myred}{\textbf{EXAM TRAP:}} Confidentiality = \textcolor{myblue}{opaque decision-making process}!
	\end{frame}
	
	\begin{frame}{Q63: Specialized Knowledge - Tutorial}
		\fontsize{6.5}{7.5}\selectfont
		\textbf{QUESTION: What is specialized knowledge as a form of opaqueness?}
		
		\vspace{0.1cm}
		\textbf{OPTIONS:}
		\begin{enumerate}
			\item When everyone can understand algorithms
			\item When understanding algorithms requires specialized training
			\item When algorithms are always simple
			\item When no knowledge is needed
		\end{enumerate}
		
		\vspace{0.1cm}
		\textcolor{myblue}{\textbf{ANSWER: B}}
		
		\vspace{0.1cm}
		\textbf{DETAILED EXPLANATION:}
		\begin{itemize}
			\item Even with access to source code
			\item May not be able to make sense of it
			\item Requires specialized training
			\item Not unique to AI systems
		\end{itemize}
		
		\vspace{0.1cm}
		\textcolor{myred}{\textbf{WHY OTHER OPTIONS ARE WRONG:}}
		\begin{itemize}
			\item \textcolor{myred}{A:} Specialized knowledge means NOT everyone understands
			\item \textcolor{myred}{C:} Algorithms may be complex
			\item \textcolor{myred}{D:} Knowledge IS needed
		\end{itemize}
		
		\textcolor{myred}{\textbf{EXAM TRAP:}} Specialized knowledge = \textcolor{myblue}{requires training to understand}!
	\end{frame}
	
	\begin{frame}{Q64: Reasons for Opacity - Tutorial}
		\fontsize{6.5}{7.5}\selectfont
		\textbf{QUESTION: What is a valid reason for algorithmic opacity?}
		
		\vspace{0.1cm}
		\textbf{OPTIONS:}
		\begin{enumerate}
			\item To hide discrimination
			\item Security (spam filters) and proprietary interests
			\item To avoid accountability
			\item To confuse users
		\end{enumerate}
		
		\vspace{0.1cm}
		\textcolor{myblue}{\textbf{ANSWER: B}}
		
		\vspace{0.1cm}
		\textbf{DETAILED EXPLANATION:}
		\begin{itemize}
			\item Security: making code public enables circumvention
			\item Example: Spam filters would be easier to bypass
			\item Proprietary interests: trade secrets
			\item Must balance transparency with legitimate interests
		\end{itemize}
		
		\vspace{0.1cm}
		\textcolor{myred}{\textbf{WHY OTHER OPTIONS ARE WRONG:}}
		\begin{itemize}
			\item \textcolor{myred}{A:} Hiding discrimination is NOT valid
			\item \textcolor{myred}{C:} Avoiding accountability is NOT valid
			\item \textcolor{myred}{D:} Confusing users is NOT valid
		\end{itemize}
		
		\textcolor{myred}{\textbf{EXAM TRAP:}} Legitimate reasons include \textcolor{myblue}{security and proprietary interests}!
	\end{frame}
	
	\begin{frame}{Q65: Explainability Summary - Tutorial}
		\fontsize{6.5}{7.5}\selectfont
		\textbf{QUESTION: What is the most comprehensive statement about explainability?}
		
		\vspace{0.1cm}
		\textbf{OPTIONS:}
		\begin{enumerate}
			\item Explainability is always easy to achieve
			\item Explainability, interpretability, and transparency are essential for accountability and trust in AI systems
			\item Explainability is never needed
			\item Only black-box models are used
		\end{enumerate}
		
		\vspace{0.1cm}
		\textcolor{myblue}{\textbf{ANSWER: B}}
		
		\vspace{0.1cm}
		\textbf{DETAILED EXPLANATION:}
		\begin{itemize}
			\item \textbf{Explainability:} Ex-post understanding of decisions
			\item \textbf{Interpretability:} Inherent understanding of model behavior
			\item \textbf{Transparency:} Openness about design and processes
			\item All essential for accountability and trust
		\end{itemize}
		
		\vspace{0.1cm}
		\textcolor{myred}{\textbf{WHY OTHER OPTIONS ARE WRONG:}}
		\begin{itemize}
			\item \textcolor{myred}{A:} Explainability can be challenging
			\item \textcolor{myred}{C:} Explainability IS needed for accountability
			\item \textcolor{myred}{D:} Many models are interpretable
		\end{itemize}
		
		\textcolor{myred}{\textbf{EXAM TRAP:}} All three are \textcolor{myblue}{essential for accountability and trust}!
	\end{frame}
	
	% ============================================================
	% SECTION 5: AUTONOMY AND MANIPULATION - QUESTIONS 66-70
	% ============================================================
	
	\begin{frame}{Q66: Human Autonomy Definition - Tutorial}
		\fontsize{6.5}{7.5}\selectfont
		\textbf{QUESTION: What is human autonomy in the context of AI?}
		
		\vspace{0.1cm}
		\textbf{OPTIONS:}
		\begin{enumerate}
			\item The ability to make any decision
			\item The ability to hold beliefs, make, and execute decisions of our own
			\item The ability to control AI systems
			\item The ability to avoid all technology
		\end{enumerate}
		
		\vspace{0.1cm}
		\textcolor{myblue}{\textbf{ANSWER: B}}
		
		\vspace{0.1cm}
		\textbf{DETAILED EXPLANATION:}
		\begin{itemize}
			\item Two dimensions:
			\item \textbf{Internal:} Values, beliefs, decisions that are our own
			\item \textbf{External:} Freedom and opportunity to execute decisions
			\item Not automatically lost by delegating tasks
			\item Requires meaningful human oversight
		\end{itemize}
		
		\vspace{0.1cm}
		\textcolor{myred}{\textbf{WHY OTHER OPTIONS ARE WRONG:}}
		\begin{itemize}
			\item \textcolor{myred}{A:} Autonomy is more than making any decision
			\item \textcolor{myred}{C:} This is control, not autonomy
			\item \textcolor{myred}{D:} Avoiding technology is not autonomy
		\end{itemize}
		
		\textcolor{myred}{\textbf{EXAM TRAP:}} Autonomy = \textcolor{myblue}{own beliefs and decisions + ability to execute}!
	\end{frame}
	
	\begin{frame}{Q67: Cambridge Analytica Scandal - Tutorial}
		\fontsize{6.5}{7.5}\selectfont
		\textbf{QUESTION: What did the Cambridge Analytica scandal demonstrate?}
		
		\vspace{0.1cm}
		\textbf{OPTIONS:}
		\begin{enumerate}
			\item The benefits of data collection
			\item How big data could be used to manipulate voter behavior
			\item That data privacy is not important
			\item That AI never affects elections
		\end{enumerate}
		
		\vspace{0.1cm}
		\textcolor{myblue}{\textbf{ANSWER: B}}
		
		\vspace{0.1cm}
		\textbf{DETAILED EXPLANATION:}
		\begin{itemize}
			\item Collected Facebook user data
			\item Inferred personality profiles from online behavior
			\item Targeted political advertising
			\item Demonstrated potential for manipulation
		\end{itemize}
		
		\vspace{0.1cm}
		\textcolor{myred}{\textbf{WHY OTHER OPTIONS ARE WRONG:}}
		\begin{itemize}
			\item \textcolor{myred}{A:} The scandal showed MISUSE of data
			\item \textcolor{myred}{C:} The scandal highlighted privacy concerns
			\item \textcolor{myred}{D:} AI CAN affect elections
		\end{itemize}
		
		\textcolor{myred}{\textbf{EXAM TRAP:}} Cambridge Analytica = \textcolor{myblue}{political manipulation through data}!
	\end{frame}
	
	\begin{frame}{Q68: Facebook Emotional Contagion Experiment - Tutorial}
		\fontsize{6.5}{7.5}\selectfont
		\textbf{QUESTION: What did the Facebook emotional contagion experiment show?}
		
		\vspace{0.1cm}
		\textbf{OPTIONS:}
		\begin{enumerate}
			\item That emotions cannot be transmitted online
			\item That recommendation algorithms can shape and manipulate user experience
			\item That users always consent to experiments
			\item That AI has no effect on emotions
		\end{enumerate}
		
		\vspace{0.1cm}
		\textcolor{myblue}{\textbf{ANSWER: B}}
		
		\vspace{0.1cm}
		\textbf{DETAILED EXPLANATION:}
		\begin{itemize}
			\item Researchers changed users' timelines
			\item Displayed predominantly positive or negative posts
			\item Users' own posts became more positive/negative
			\item No consent from users
			\item Caused public outcry
		\end{itemize}
		
		\vspace{0.1cm}
		\textcolor{myred}{\textbf{WHY OTHER OPTIONS ARE WRONG:}}
		\begin{itemize}
			\item \textcolor{myred}{A:} Emotions CAN be transmitted online
			\item \textcolor{myred}{C:} Users did NOT consent
			\item \textcolor{myred}{D:} AI CAN affect emotions
		\end{itemize}
		
		\textcolor{myred}{\textbf{EXAM TRAP:}} Experiment showed \textcolor{myblue}{AI can manipulate emotional state}!
	\end{frame}
	
	\begin{frame}{Q69: AI Manipulation Prevention - Tutorial}
		\fontsize{6.5}{7.5}\selectfont
		\textbf{QUESTION: How can AI manipulation be prevented?}
		
		\vspace{0.1cm}
		\textbf{OPTIONS:}
		\begin{enumerate}
			\item By ignoring all risks
			\item Through human oversight and regulatory oversight
			\item By eliminating all AI
			\item By never collecting data
		\end{enumerate}
		
		\vspace{0.1cm}
		\textcolor{myblue}{\textbf{ANSWER: B}}
		
		\vspace{0.1cm}
		\textbf{DETAILED EXPLANATION:}
		\begin{itemize}
			\item Human oversight throughout development
			\item Regulatory oversight (e.g., EU AI Act)
			\item Transparency obligations
			\item Prohibiting subliminal techniques
			\item Balancing benefits vs. risks
		\end{itemize}
		
		\vspace{0.1cm}
		\textcolor{myred}{\textbf{WHY OTHER OPTIONS ARE WRONG:}}
		\begin{itemize}
			\item \textcolor{myred}{A:} Ignoring risks is not prevention
			\item \textcolor{myred}{C:} Eliminating AI is unrealistic
			\item \textcolor{myred}{D:} Some data collection is necessary
		\end{itemize}
		
		\textcolor{myred}{\textbf{EXAM TRAP:}} Prevention requires \textcolor{myblue}{human and regulatory oversight}!
	\end{frame}
	
	\begin{frame}{Q70: Autonomy and Manipulation Summary - Tutorial}
		\fontsize{6.5}{7.5}\selectfont
		\textbf{QUESTION: What is the most comprehensive statement about autonomy and manipulation?}
		
		\vspace{0.1cm}
		\textbf{OPTIONS:}
		\begin{enumerate}
			\item AI never affects autonomy
			\item AI systems can shape our experience and potentially manipulate us, requiring oversight and regulation
			\item Manipulation is always beneficial
			\item Autonomy is not important
		\end{enumerate}
		
		\vspace{0.1cm}
		\textcolor{myblue}{\textbf{ANSWER: B}}
		
		\vspace{0.1cm}
		\textbf{DETAILED EXPLANATION:}
		\begin{itemize}
			\item AI shapes online experience
			\item Can influence behavior through recommendation algorithms
			\item Examples: Cambridge Analytica, emotional contagion
			\item Requires oversight and regulation
		\end{itemize}
		
		\vspace{0.1cm}
		\textcolor{myred}{\textbf{WHY OTHER OPTIONS ARE WRONG:}}
		\begin{itemize}
			\item \textcolor{myred}{A:} AI CAN affect autonomy
			\item \textcolor{myred}{C:} Manipulation is generally harmful
			\item \textcolor{myred}{D:} Autonomy IS important
		\end{itemize}
		
		\textcolor{myred}{\textbf{EXAM TRAP:}} AI can \textcolor{myblue}{shape experience and manipulate} - requires oversight!
	\end{frame}
	
	% ============================================================
	% SECTION 6: SAFETY AND WELL-BEING - QUESTIONS 71-75
	% ============================================================
	
	\begin{frame}{Q71: Uber Crash Example - Tutorial}
		\fontsize{6.5}{7.5}\selectfont
		\textbf{QUESTION: What did the Uber driverless vehicle crash demonstrate?}
		
		\vspace{0.1cm}
		\textbf{OPTIONS:}
		\begin{enumerate}
			\item That autonomous vehicles are always safe
			\item The risks of over-dependence on AI in critical systems
			\item That safety mechanisms are never needed
			\item That AI never fails
		\end{enumerate}
		
		\vspace{0.1cm}
		\textcolor{myblue}{\textbf{ANSWER: B}}
		
		\vspace{0.1cm}
		\textbf{DETAILED EXPLANATION:}
		\begin{itemize}
			\item Vehicle's systems failed to recognize pedestrian
			\item Safety mechanisms were deactivated
			\item Human safety driver was distracted
			\item Result of automation bias
			\item Need for robust testing and validation
		\end{itemize}
		
		\vspace{0.1cm}
		\textcolor{myred}{\textbf{WHY OTHER OPTIONS ARE WRONG:}}
		\begin{itemize}
			\item \textcolor{myred}{A:} The crash showed they are NOT always safe
			\item \textcolor{myred}{C:} Safety mechanisms ARE needed
			\item \textcolor{myred}{D:} AI CAN fail
		\end{itemize}
		
		\textcolor{myred}{\textbf{EXAM TRAP:}} Uber crash = \textcolor{myblue}{risks of over-reliance on AI}!
	\end{frame}
	
	\begin{frame}{Q72: Automation Bias Definition - Tutorial}
		\fontsize{6.5}{7.5}\selectfont
		\textbf{QUESTION: What is automation bias?}
		
		\vspace{0.1cm}
		\textbf{OPTIONS:}
		\begin{enumerate}
			\item The tendency to distrust automated systems
			\item The tendency of humans to over-rely on automated systems
			\item The tendency to always use manual methods
			\item The tendency to ignore technology
		\end{enumerate}
		
		\vspace{0.1cm}
		\textcolor{myblue}{\textbf{ANSWER: B}}
		
		\vspace{0.1cm}
		\textbf{DETAILED EXPLANATION:}
		\begin{itemize}
			\item Over-reliance on automation
			\item Leads to complacency and reduced vigilance
			\item Errors of oversight
			\item Particularly problematic in critical situations
		\end{itemize}
		
		\vspace{0.1cm}
		\textcolor{myred}{\textbf{WHY OTHER OPTIONS ARE WRONG:}}
		\begin{itemize}
			\item \textcolor{myred}{A:} This is distrust, not over-reliance
			\item \textcolor{myred}{C:} Automation bias is about using automation
			\item \textcolor{myred}{D:} Automation bias is about over-relying
		\end{itemize}
		
		\textcolor{myred}{\textbf{EXAM TRAP:}} Automation bias = \textcolor{myblue}{over-reliance on automated systems}!
	\end{frame}
	
	\begin{frame}{Q73: Skill Atrophy - Tutorial}
		\fontsize{6.5}{7.5}\selectfont
		\textbf{QUESTION: What is skill atrophy in the context of AI?}
		
		\vspace{0.1cm}
		\textbf{OPTIONS:}
		\begin{enumerate}
			\item Improving skills through AI use
			\item Decay in essential skills from prolonged reliance on automation
			\item Learning new skills quickly
			\item Maintaining all skills
		\end{enumerate}
		
		\vspace{0.1cm}
		\textcolor{myblue}{\textbf{ANSWER: B}}
		
		\vspace{0.1cm}
		\textbf{DETAILED EXPLANATION:}
		\begin{itemize}
			\item Operators become passive observers
			\item Move from active decision makers
			\item Essential skills deteriorate
			\item Result of automation bias
		\end{itemize}
		
		\vspace{0.1cm}
		\textcolor{myred}{\textbf{WHY OTHER OPTIONS ARE WRONG:}}
		\begin{itemize}
			\item \textcolor{myred}{A:} This is skill improvement
			\item \textcolor{myred}{C:} Skill atrophy is about DECAY, not learning
			\item \textcolor{myred}{D:} Skill atrophy is about LOSING skills
		\end{itemize}
		
		\textcolor{myred}{\textbf{EXAM TRAP:}} Skill atrophy = \textcolor{myblue}{decay from over-reliance on automation}!
	\end{frame}
	
	\begin{frame}{Q74: Automation Bias Mitigation - Tutorial}
		\fontsize{6.5}{7.5}\selectfont
		\textbf{QUESTION: How can automation bias be mitigated?}
		
		\vspace{0.1cm}
		\textbf{OPTIONS:}
		\begin{enumerate}
			\item By removing all automation
			\item By informing users and requiring periodic human input
			\item By ignoring the problem
			\item By increasing automation
		\end{enumerate}
		
		\vspace{0.1cm}
		\textcolor{myblue}{\textbf{ANSWER: B}}
		
		\vspace{0.1cm}
		\textbf{DETAILED EXPLANATION:}
		\begin{itemize}
			\item Inform users/operators of the risks
			\item Ensure understanding of limitations
			\item Implement design mechanisms requiring periodic human input
			\item Example: Eye tracking technology with alerts
		\end{itemize}
		
		\vspace{0.1cm}
		\textcolor{myred}{\textbf{WHY OTHER OPTIONS ARE WRONG:}}
		\begin{itemize}
			\item \textcolor{myred}{A:} Removing all automation is unrealistic
			\item \textcolor{myred}{C:} Ignoring the problem is not mitigation
			\item \textcolor{myred}{D:} Increasing automation worsens the problem
		\end{itemize}
		
		\textcolor{myred}{\textbf{EXAM TRAP:}} Mitigation = \textcolor{myblue}{user education + periodic human verification}!
	\end{frame}
	
	\begin{frame}{Q75: Mental Health AI Risks - Tutorial}
		\fontsize{6.5}{7.5}\selectfont
		\textbf{QUESTION: What is a concern with AI-driven mental health support?}
		
		\vspace{0.1cm}
		\textbf{OPTIONS:}
		\begin{enumerate}
			\item AI systems have perfect empathy
			\item AI systems have no empathy, which is crucial for mental health support
			\item AI systems are always better than humans
			\item AI systems never fail in mental health
		\end{enumerate}
		
		\vspace{0.1cm}
		\textcolor{myblue}{\textbf{ANSWER: B}}
		
		\vspace{0.1cm}
		\textbf{DETAILED EXPLANATION:}
		\begin{itemize}
			\item AI can give seemingly empathetic responses
			\item But AI has no actual empathy
			\item Positive effects dwindle when patients know it's AI-generated
			\item Some issues should remain in human hands
		\end{itemize}
		
		\vspace{0.1cm}
		\textcolor{myred}{\textbf{WHY OTHER OPTIONS ARE WRONG:}}
		\begin{itemize}
			\item \textcolor{myred}{A:} AI does NOT have perfect empathy
			\item \textcolor{myred}{C:} AI is NOT always better than humans
			\item \textcolor{myred}{D:} AI CAN fail in mental health
		\end{itemize}
		
		\textcolor{myred}{\textbf{EXAM TRAP:}} AI lacks \textcolor{myblue}{genuine empathy} - critical for mental health!
	\end{frame}
	
	% ============================================================
	% SECTION 7: REPUTATIONAL RISKS - QUESTIONS 76-80
	% ============================================================
	
	\begin{frame}{Q76: Reputational Risk Causes - Tutorial}
		\fontsize{6.5}{7.5}\selectfont
		\textbf{QUESTION: What are the three main causes of AI-related reputational damage?}
		
		\vspace{0.1cm}
		\textbf{OPTIONS:}
		\begin{enumerate}
			\item Speed, cost, and accuracy
			\item Privacy breaches, algorithmic bias, and lack of explainability
			\item Size, complexity, and speed
			\item Training data, model type, and deployment
		\end{enumerate}
		
		\vspace{0.1cm}
		\textcolor{myblue}{\textbf{ANSWER: B}}
		
		\vspace{0.1cm}
		\textbf{DETAILED EXPLANATION:}
		\begin{itemize}
			\item \textbf{Privacy breaches:} Cambridge Analytica example
			\item \textbf{Algorithmic bias:} Dutch child benefits scandal
			\item \textbf{Lack of explainability:} Customers denied loans without reasons
		\end{itemize}
		
		\vspace{0.1cm}
		\textcolor{myred}{\textbf{WHY OTHER OPTIONS ARE WRONG:}}
		\begin{itemize}
			\item \textcolor{myred}{A:} Speed, cost, accuracy are not the main causes
			\item \textcolor{myred}{C:} Size, complexity, speed are not the main causes
			\item \textcolor{myred}{D:} Training data, model type, deployment are not the main causes
		\end{itemize}
		
		\textcolor{myred}{\textbf{EXAM TRAP:}} Three main causes = \textcolor{myblue}{privacy, bias, explainability}!
	\end{frame}
	
	\begin{frame}{Q77: Dutch Child Benefits Scandal - Tutorial}
		\fontsize{6.5}{7.5}\selectfont
		\textbf{QUESTION: What happened in the Dutch child benefits scandal?}
		
		\vspace{0.1cm}
		\textbf{OPTIONS:}
		\begin{enumerate}
			\item An algorithm correctly identified fraud
			\item A fraud-detection algorithm falsely denied benefits to thousands of citizens
			\item All claims were approved
			\item The algorithm worked perfectly
		\end{enumerate}
		
		\vspace{0.1cm}
		\textcolor{myblue}{\textbf{ANSWER: B}}
		
		\vspace{0.1cm}
		\textbf{DETAILED EXPLANATION:}
		\begin{itemize}
			\item Algorithm deployed by Dutch government
			\item Falsely denied and reclaimed child benefits
			\item Thousands of citizens affected
			\item Dire consequences for families
			\item Dutch prime minister resigned
		\end{itemize}
		
		\vspace{0.1cm}
		\textcolor{myred}{\textbf{WHY OTHER OPTIONS ARE WRONG:}}
		\begin{itemize}
			\item \textcolor{myred}{A:} The algorithm incorrectly identified fraud
			\item \textcolor{myred}{C:} Many claims were DENIED
			\item \textcolor{myred}{D:} The algorithm failed catastrophically
		\end{itemize}
		
		\textcolor{myred}{\textbf{EXAM TRAP:}} Dutch scandal = \textcolor{myblue}{algorithmic bias with severe consequences}!
	\end{frame}
	
	\begin{frame}{Q78: Competence Criticism - Tutorial}
		\fontsize{6.5}{7.5}\selectfont
		\textbf{QUESTION: What is competence criticism of companies?}
		
		\vspace{0.1cm}
		\textbf{OPTIONS:}
		\begin{enumerate}
			\item Criticism that the company lacks relevant expertise
			\item Criticism that the company has misaligned values
			\item Criticism that the company is too successful
			\item Criticism that the company is too transparent
		\end{enumerate}
		
		\vspace{0.1cm}
		\textcolor{myblue}{\textbf{ANSWER: A}}
		
		\vspace{0.1cm}
		\textbf{DETAILED EXPLANATION:}
		\begin{itemize}
			\item Perception that company lacks expertise
			\item Insufficient capacities for safe/fair algorithms
			\item Technical failings, unfairness, or lack of explainability
			\item Data not stored safely
		\end{itemize}
		
		\vspace{0.1cm}
		\textcolor{myred}{\textbf{WHY OTHER OPTIONS ARE WRONG:}}
		\begin{itemize}
			\item \textcolor{myred}{B:} This is value alignment criticism
			\item \textcolor{myred}{C:} This is not criticism
			\item \textcolor{myred}{D:} Transparency is generally positive
		\end{itemize}
		
		\textcolor{myred}{\textbf{EXAM TRAP:}} Competence = \textcolor{myblue}{lack of expertise or capacity}!
	\end{frame}
	
	\begin{frame}{Q79: Value Alignment Criticism - Tutorial}
		\fontsize{6.5}{7.5}\selectfont
		\textbf{QUESTION: What is value alignment criticism of companies?}
		
		\vspace{0.1cm}
		\textbf{OPTIONS:}
		\begin{enumerate}
			\item Criticism that the company lacks relevant expertise
			\item Criticism that the company has misaligned values or acted despite knowing risks
			\item Criticism that the company is too successful
			\item Criticism that the company is too transparent
		\end{enumerate}
		
		\vspace{0.1cm}
		\textcolor{myblue}{\textbf{ANSWER: B}}
		
		\vspace{0.1cm}
		\textbf{DETAILED EXPLANATION:}
		\begin{itemize}
			\item Company fully aware of risks but proceeded
			\item Example: Facebook emotional contagion experiment
			\item Lack of due diligence
			\item Violation of user trust and expectations
		\end{itemize}
		
		\vspace{0.1cm}
		\textcolor{myred}{\textbf{WHY OTHER OPTIONS ARE WRONG:}}
		\begin{itemize}
			\item \textcolor{myred}{A:} This is competence criticism
			\item \textcolor{myred}{C:} This is not criticism
			\item \textcolor{myred}{D:} Transparency is generally positive
		\end{itemize}
		
		\textcolor{myred}{\textbf{EXAM TRAP:}} Value alignment = \textcolor{myblue}{misaligned values or knowingly proceeding with risks}!
	\end{frame}
	
	\begin{frame}{Q80: Reputational Risk Mitigation - Tutorial}
		\fontsize{6.5}{7.5}\selectfont
		\textbf{QUESTION: What is NOT a reputational risk mitigation strategy?}
		
		\vspace{0.1cm}
		\textbf{OPTIONS:}
		\begin{enumerate}
			\item Continuous monitoring and evaluation
			\item Ignoring incidents and hoping they go away
			\item Ethical AI frameworks
			\item Stakeholder engagement and communication
		\end{enumerate}
		
		\vspace{0.1cm}
		\textcolor{myblue}{\textbf{ANSWER: B}}
		
		\vspace{0.1cm}
		\textbf{DETAILED EXPLANATION:}
		\begin{itemize}
			\item \textbf{Monitoring:} Regular risk assessment
			\item \textbf{Transparency:} Explainable algorithms, user awareness
			\item \textbf{Ethical frameworks:} Guidelines and oversight
			\item \textbf{Stakeholder engagement:} Communication and trust-building
			\item \textbf{Crisis management:} Plan for AI-related incidents
		\end{itemize}
		
		\vspace{0.1cm}
		\textcolor{myred}{\textbf{WHY OTHER OPTIONS ARE WRONG:}}
		\begin{itemize}
			\item \textcolor{myred}{A:} This IS a mitigation strategy
			\item \textcolor{myred}{C:} This IS a mitigation strategy
			\item \textcolor{myred}{D:} This IS a mitigation strategy
		\end{itemize}
		
		\textcolor{myred}{\textbf{EXAM TRAP:}} Ignoring incidents is \textcolor{myblue}{NOT a valid mitigation strategy}!
	\end{frame}
	
	% ============================================================
	% SECTION 8: EXISTENTIAL AND GLOBAL RISKS - QUESTIONS 81-85
	% ============================================================
	
	\begin{frame}{Q81: Superintelligence Definition - Tutorial}
		\fontsize{6.5}{7.5}\selectfont
		\textbf{QUESTION: What is superintelligence?}
		
		\vspace{0.1cm}
		\textbf{OPTIONS:}
		\begin{enumerate}
			\item An AI system that is faster than humans
			\item An AI system that operates beyond human-level intelligence
			\item An AI system that is less intelligent than humans
			\item An AI system that only performs one task
		\end{enumerate}
		
		\vspace{0.1cm}
		\textcolor{myblue}{\textbf{ANSWER: B}}
		
		\vspace{0.1cm}
		\textbf{DETAILED EXPLANATION:}
		\begin{itemize}
			\item AI beyond human-level intelligence
			\item Could act in ways threatening to humanity
			\item Concern: well-intentioned AI with misaligned values
			\item Central to existential risk discussions
		\end{itemize}
		
		\vspace{0.1cm}
		\textcolor{myred}{\textbf{WHY OTHER OPTIONS ARE WRONG:}}
		\begin{itemize}
			\item \textcolor{myred}{A:} Speed is different from intelligence level
			\item \textcolor{myred}{C:} Superintelligence is ABOVE human level
			\item \textcolor{myred}{D:} Superintelligence would be general, not narrow
		\end{itemize}
		
		\textcolor{myred}{\textbf{EXAM TRAP:}} Superintelligence = \textcolor{myblue}{beyond human-level intelligence}!
	\end{frame}
	
	\begin{frame}{Q82: AI Alignment - Tutorial}
		\fontsize{6.5}{7.5}\selectfont
		\textbf{QUESTION: What is AI alignment?}
		
		\vspace{0.1cm}
		\textbf{OPTIONS:}
		\begin{enumerate}
			\item Aligning AI systems with human values
			\item Aligning AI systems with maximum efficiency
			\item Aligning AI systems with profit maximization
			\item Aligning AI systems with speed
		\end{enumerate}
		
		\vspace{0.1cm}
		\textcolor{myblue}{\textbf{ANSWER: A}}
		
		\vspace{0.1cm}
		\textbf{DETAILED EXPLANATION:}
		\begin{itemize}
			\item Ensuring AI systems' means and goals align with human values
			\item Example: Paperclip maximizer scenario
			\item Misaligned AI could cause catastrophic harm
			\item Central concern for existential risk
		\end{itemize}
		
		\vspace{0.1cm}
		\textcolor{myred}{\textbf{WHY OTHER OPTIONS ARE WRONG:}}
		\begin{itemize}
			\item \textcolor{myred}{B:} Efficiency is not alignment
			\item \textcolor{myred}{C:} Profit is not alignment
			\item \textcolor{myred}{D:} Speed is not alignment
		\end{itemize}
		
		\textcolor{myred}{\textbf{EXAM TRAP:}} AI alignment = \textcolor{myblue}{matching AI to human values}!
	\end{frame}
	
	\begin{frame}{Q83: AI Arms Race - Tutorial}
		\fontsize{6.5}{7.5}\selectfont
		\textbf{QUESTION: What is the AI arms race concern?}
		
		\vspace{0.1cm}
		\textbf{OPTIONS:}
		\begin{enumerate}
			\item Competing nations developing AI for peaceful purposes
			\item Risks from different stakeholders creating powerful AI without adequate safety measures
			\item A race to create the fastest AI
			\item A race to create the most profitable AI
		\end{enumerate}
		
		\vspace{0.1cm}
		\textcolor{myblue}{\textbf{ANSWER: B}}
		
		\vspace{0.1cm}
		\textbf{DETAILED EXPLANATION:}
		\begin{itemize}
			\item Multiple stakeholders creating powerful AI
			\item Without adequate safety measures
			\item Could lead to catastrophic conflicts
			\item Part of existential risk discussions
		\end{itemize}
		
		\vspace{0.1cm}
		\textcolor{myred}{\textbf{WHY OTHER OPTIONS ARE WRONG:}}
		\begin{itemize}
			\item \textcolor{myred}{A:} The concern is about UNSAFE development
			\item \textcolor{myred}{C:} Speed is not the primary concern
			\item \textcolor{myred}{D:} Profit is not the primary concern
		\end{itemize}
		
		\textcolor{myred}{\textbf{EXAM TRAP:}} AI arms race = \textcolor{myblue}{unsafe AI development competition}!
	\end{frame}
	
	\begin{frame}{Q84: Economic Risks from AI - Tutorial}
		\fontsize{6.5}{7.5}\selectfont
		\textbf{QUESTION: What is a major economic risk from AI?}
		
		\vspace{0.1cm}
		\textbf{OPTIONS:}
		\begin{enumerate}
			\item AI always creates more jobs
			\item AI never affects employment
			\item Displacement of jobs by AI and widening inequality
			\item AI has no economic impact
		\end{enumerate}
		
		\vspace{0.1cm}
		\textcolor{myblue}{\textbf{ANSWER: C}}
		
		\vspace{0.1cm}
		\textbf{DETAILED EXPLANATION:}
		\begin{itemize}
			\item Jobs could be displaced by automation
			\item High-skill jobs also at risk (accounting, legal)
			\item Widening gap between AI-skilled and unskilled
			\item Need for workforce re-skilling
			\item Predictions vary (Goldman Sachs: 1/4 jobs could be automated)
		\end{itemize}
		
		\vspace{0.1cm}
		\textcolor{myred}{\textbf{WHY OTHER OPTIONS ARE WRONG:}}
		\begin{itemize}
			\item \textcolor{myred}{A:} AI may DESTROY some jobs
			\item \textcolor{myred}{B:} AI DOES affect employment
			\item \textcolor{myred}{D:} AI HAS significant economic impact
		\end{itemize}
		
		\textcolor{myred}{\textbf{EXAM TRAP:}} Economic risks = \textcolor{myblue}{job displacement and inequality}!
	\end{frame}
	
	\begin{frame}{Q85: Global Inequality from AI - Tutorial}
		\fontsize{6.5}{7.5}\selectfont
		\textbf{QUESTION: How could AI exacerbate global inequality?}
		
		\vspace{0.1cm}
		\textbf{OPTIONS:}
		\begin{enumerate}
			\item AI is evenly distributed globally
			\item Nations at the center of AI development leap ahead, leaving others behind
			\item AI has no effect on global inequality
			\item AI reduces inequality automatically
		\end{enumerate}
		
		\vspace{0.1cm}
		\textcolor{myblue}{\textbf{ANSWER: B}}
		
		\vspace{0.1cm}
		\textbf{DETAILED EXPLANATION:}
		\begin{itemize}
			\item Most large AI firms based in the US
			\item Technological adoption varies widely
			\item AI advancement could widen global inequality
			\item Nations in AI development leap ahead
			\item Others left behind in economic growth and influence
		\end{itemize}
		
		\vspace{0.1cm}
		\textcolor{myred}{\textbf{WHY OTHER OPTIONS ARE WRONG:}}
		\begin{itemize}
			\item \textcolor{myred}{A:} AI is NOT evenly distributed
			\item \textcolor{myred}{C:} AI CAN exacerbate inequality
			\item \textcolor{myred}{D:} AI does NOT automatically reduce inequality
		\end{itemize}
		
		\textcolor{myred}{\textbf{EXAM TRAP:}} AI could \textcolor{myblue}{widen global inequality}!
	\end{frame}
	
	% ============================================================
	% SECTION 9: GENAI RISKS - QUESTIONS 86-100
	% ============================================================
	
	\begin{frame}{Q86: GenAI Adoption Statistics - Tutorial}
		\fontsize{6.5}{7.5}\selectfont
		\textbf{QUESTION: According to McKinsey (2025), what percentage of organizations use GenAI?}
		
		\vspace{0.1cm}
		\textbf{OPTIONS:}
		\begin{enumerate}
			\item 33\%
			\item 71\%
			\item 50\%
			\item 95\%
		\end{enumerate}
		
		\vspace{0.1cm}
		\textcolor{myblue}{\textbf{ANSWER: B}}
		
		\vspace{0.1cm}
		\textbf{DETAILED EXPLANATION:}
		\begin{itemize}
			\item 71\% of organizations regularly use GenAI
			\item Up from 33\% in 2023
			\item Particularly high in digitized industries
			\item Banking and financial services: 58\% embrace GenAI
		\end{itemize}
		
		\vspace{0.1cm}
		\textcolor{myred}{\textbf{WHY OTHER OPTIONS ARE WRONG:}}
		\begin{itemize}
			\item \textcolor{myred}{A:} This was the 2023 figure
			\item \textcolor{myred}{C:} The figure is higher
			\item \textcolor{myred}{D:} This is the BCG figure for professional services
		\end{itemize}
		
		\textcolor{myred}{\textbf{EXAM TRAP:}} GenAI adoption = \textcolor{myblue}{71\% of organizations} (2025 McKinsey)!
	\end{frame}
	
	\begin{frame}{Q87: Hallucination Definition - Tutorial}
		\fontsize{6.5}{7.5}\selectfont
		\textbf{QUESTION: What is hallucination in GenAI?}
		
		\vspace{0.1cm}
		\textbf{OPTIONS:}
		\begin{enumerate}
			\item Always correct outputs
			\item Factually incorrect, misleading, or fabricated outputs that appear credible
			\item Outputs that are never credible
			\item Outputs that are always correct
		\end{enumerate}
		
		\vspace{0.1cm}
		\textcolor{myblue}{\textbf{ANSWER: B}}
		
		\vspace{0.1cm}
		\textbf{DETAILED EXPLANATION:}
		\begin{itemize}
			\item Generate factually incorrect output
			\item Appear credible to non-experts
			\item Problematic in legal, financial, healthcare contexts
			\item Can lead to reputational damage and legal liability
		\end{itemize}
		
		\vspace{0.1cm}
		\textcolor{myred}{\textbf{WHY OTHER OPTIONS ARE WRONG:}}
		\begin{itemize}
			\item \textcolor{myred}{A:} Hallucinations are INCORRECT
			\item \textcolor{myred}{C:} Hallucinations APPEAR credible
			\item \textcolor{myred}{D:} Hallucinations are NOT correct
		\end{itemize}
		
		\textcolor{myred}{\textbf{EXAM TRAP:}} Hallucination = \textcolor{myblue}{incorrect but credible outputs}!
	\end{frame}
	
	\begin{frame}{Q88: Air Canada Chatbot Case - Tutorial}
		\fontsize{6.5}{7.5}\selectfont
		\textbf{QUESTION: What happened in the Air Canada chatbot case?}
		
		\vspace{0.1cm}
		\textbf{OPTIONS:}
		\begin{enumerate}
			\item The chatbot worked perfectly
			\item The chatbot offered a discount that shouldn't have been available, and the company was held responsible
			\item The chatbot was a separate legal entity
			\item The company won the case
		\end{enumerate}
		
		\vspace{0.1cm}
		\textcolor{myblue}{\textbf{ANSWER: B}}
		
		\vspace{0.1cm}
		\textbf{DETAILED EXPLANATION:}
		\begin{itemize}
			\item Chatbot offered a passenger a discount
			\item Discount shouldn't have been available
			\item Air Canada argued chatbot was separate legal entity
			\item Tribunal rejected argument
			\item Company responsible for chatbot outputs
			\item Awarded damages and tribunal fees
		\end{itemize}
		
		\vspace{0.1cm}
		\textcolor{myred}{\textbf{WHY OTHER OPTIONS ARE WRONG:}}
		\begin{itemize}
			\item \textcolor{myred}{A:} The chatbot made an error
			\item \textcolor{myred}{C:} The argument was REJECTED
			\item \textcolor{myred}{D:} Air Canada LOST the case
		\end{itemize}
		
		\textcolor{myred}{\textbf{EXAM TRAP:}} Companies are \textcolor{myblue}{liable for chatbot outputs}!
	\end{frame}
	
	\begin{frame}{Q89: Unpredictability in GenAI - Tutorial}
		\fontsize{6.5}{7.5}\selectfont
		\textbf{QUESTION: Why is unpredictability a risk in GenAI?}
		
		\vspace{0.1cm}
		\textbf{OPTIONS:}
		\begin{enumerate}
			\item GenAI outputs are always predictable
			\item Unpredictability undermines standardization across users and time
			\item Unpredictability is always beneficial
			\item Unpredictability never affects compliance
		\end{enumerate}
		
		\vspace{0.1cm}
		\textcolor{myblue}{\textbf{ANSWER: B}}
		
		\vspace{0.1cm}
		\textbf{DETAILED EXPLANATION:}
		\begin{itemize}
			\item GenAI systems are unpredictable in output
			\item Challenges standardization across users and time
			\item Problems for auditability
			\item Affects financial advice, claims handling, legal contexts
			\item Poses compliance and reputational risks
		\end{itemize}
		
		\vspace{0.1cm}
		\textcolor{myred}{\textbf{WHY OTHER OPTIONS ARE WRONG:}}
		\begin{itemize}
			\item \textcolor{myred}{A:} GenAI outputs are NOT always predictable
			\item \textcolor{myred}{C:} Unpredictability is generally problematic
			\item \textcolor{myred}{D:} Unpredictability CAN affect compliance
		\end{itemize}
		
		\textcolor{myred}{\textbf{EXAM TRAP:}} Unpredictability undermines \textcolor{myblue}{standardization and auditability}!
	\end{frame}
	
	\begin{frame}{Q90: Prompt Injection - Tutorial}
		\fontsize{6.5}{7.5}\selectfont
		\textbf{QUESTION: What is prompt injection?}
		
		\vspace{0.1cm}
		\textbf{OPTIONS:}
		\begin{enumerate}
			\item A technique to improve outputs
			\item Hidden instructions within user inputs to cause unintended behavior
			\item A method to increase accuracy
			\item A technique to remove bias
		\end{enumerate}
		
		\vspace{0.1cm}
		\textcolor{myblue}{\textbf{ANSWER: B}}
		
		\vspace{0.1cm}
		\textbf{DETAILED EXPLANATION:}
		\begin{itemize}
			\item Hidden instructions in user inputs
			\item Cause unintended behavior
			\item One of several adversarial input techniques
			\item Can cause GenAI systems to behave unsafely
		\end{itemize}
		
		\vspace{0.1cm}
		\textcolor{myred}{\textbf{WHY OTHER OPTIONS ARE WRONG:}}
		\begin{itemize}
			\item \textcolor{myred}{A:} Prompt injection causes PROBLEMS
			\item \textcolor{myred}{C:} Prompt injection is about manipulation
			\item \textcolor{myred}{D:} Prompt injection doesn't remove bias
		\end{itemize}
		
		\textcolor{myred}{\textbf{EXAM TRAP:}} Prompt injection = \textcolor{myblue}{hidden instructions to cause unintended behavior}!
	\end{frame}
	
	\begin{frame}{Q91: Jailbreaking Definition - Tutorial}
		\fontsize{6.5}{7.5}\selectfont
		\textbf{QUESTION: What is jailbreaking in the context of GenAI?}
		
		\vspace{0.1cm}
		\textbf{OPTIONS:}
		\begin{enumerate}
			\item Making the model faster
			\item Bypassing a model's safety mechanisms through crafted prompts
			\item A method to improve outputs
			\item A technique to increase accuracy
		\end{enumerate}
		
		\vspace{0.1cm}
		\textcolor{myblue}{\textbf{ANSWER: B}}
		
		\vspace{0.1cm}
		\textbf{DETAILED EXPLANATION:}
		\begin{itemize}
			\item Crafted prompts to bypass safety mechanisms
			\item Override safety or policy constraints
			\item Can cause unintended or unsafe behavior
			\item Type of adversarial input
		\end{itemize}
		
		\vspace{0.1cm}
		\textcolor{myred}{\textbf{WHY OTHER OPTIONS ARE WRONG:}}
		\begin{itemize}
			\item \textcolor{myred}{A:} Speed is not the purpose
			\item \textcolor{myred}{C:} Jailbreaking is about bypassing safety
			\item \textcolor{myred}{D:} Jailbreaking does NOT increase accuracy
		\end{itemize}
		
		\textcolor{myred}{\textbf{EXAM TRAP:}} Jailbreaking = \textcolor{myblue}{bypassing safety mechanisms}!
	\end{frame}
	
	\begin{frame}{Q92: Data Privacy Violations in GenAI - Tutorial}
		\fontsize{6.5}{7.5}\selectfont
		\textbf{QUESTION: How can GenAI cause data privacy violations?}
		
		\vspace{0.1cm}
		\textbf{OPTIONS:}
		\begin{enumerate}
			\item GenAI never accesses personal data
			\item Prompts containing PII may be inadvertently stored, shared, or used for training
			\item Data privacy is not a concern
			\item GenAI always protects privacy
		\end{enumerate}
		
		\vspace{0.1cm}
		\textcolor{myblue}{\textbf{ANSWER: B}}
		
		\vspace{0.1cm}
		\textbf{DETAILED EXPLANATION:}
		\begin{itemize}
			\item User prompts may contain personally identifiable information (PII)
			\item Data may be inadvertently stored or shared
			\item May be used for model training without user knowledge
			\item Example: Researchers extracted emails, phone numbers from LLM
			\item FTC investigation into OpenAI for privacy issues
		\end{itemize}
		
		\vspace{0.1cm}
		\textcolor{myred}{\textbf{WHY OTHER OPTIONS ARE WRONG:}}
		\begin{itemize}
			\item \textcolor{myred}{A:} GenAI CAN access personal data
			\item \textcolor{myred}{C:} Data privacy IS a concern
			\item \textcolor{myred}{D:} GenAI does NOT always protect privacy
		\end{itemize}
		
		\textcolor{myred}{\textbf{EXAM TRAP:}} GenAI prompts may \textcolor{myblue}{inadvertently expose PII}!
	\end{frame}
	
	\begin{frame}{Q93: Copyright Issues in GenAI - Tutorial}
		\fontsize{6.5}{7.5}\selectfont
		\textbf{QUESTION: What copyright concern exists with GenAI?}
		
		\vspace{0.1cm}
		\textbf{OPTIONS:}
		\begin{enumerate}
			\item GenAI never uses copyrighted material
			\item GenAI models are trained on web-scraped content including copyrighted material
			\item Copyright laws don't apply to AI
			\item All GenAI content is public domain
		\end{enumerate}
		
		\vspace{0.1cm}
		\textcolor{myblue}{\textbf{ANSWER: B}}
		
		\vspace{0.1cm}
		\textbf{DETAILED EXPLANATION:}
		\begin{itemize}
			\item Trained on vast web-scraped content
			\item Includes copyrighted books, images, articles
			\item Lively debate and legal proceedings
			\item Example: Getty Images vs. Stability AI
			\item Using generated content may carry copyright risks
		\end{itemize}
		
		\vspace{0.1cm}
		\textcolor{myred}{\textbf{WHY OTHER OPTIONS ARE WRONG:}}
		\begin{itemize}
			\item \textcolor{myred}{A:} GenAI DOES use copyrighted material
			\item \textcolor{myred}{C:} Copyright laws DO apply
			\item \textcolor{myred}{D:} GenAI content is NOT public domain
		\end{itemize}
		
		\textcolor{myred}{\textbf{EXAM TRAP:}} GenAI training may \textcolor{myblue}{infringe copyright}!
	\end{frame}
	
	\begin{frame}{Q94: Responsibility Gaps - Tutorial}
		\fontsize{6.5}{7.5}\selectfont
		\textbf{QUESTION: What are responsibility gaps in GenAI adoption?}
		
		\vspace{0.1cm}
		\textbf{OPTIONS:}
		\begin{enumerate}
			\item Clear lines of accountability always exist
			\item GenAI tools may be adopted across teams without clear responsibility for oversight
			\item Responsibility is never an issue
			\item All employees are always responsible
		\end{enumerate}
		
		\vspace{0.1cm}
		\textcolor{myblue}{\textbf{ANSWER: B}}
		
		\vspace{0.1cm}
		\textbf{DETAILED EXPLANATION:}
		\begin{itemize}
			\item GenAI tools adopted by everyone
			\item Non-technical teams may use LLMs
			\item Example: HR team screening complaints
			\item Biased or unfair outputs possible
			\item Need clear lines of responsibility
		\end{itemize}
		
		\vspace{0.1cm}
		\textcolor{myred}{\textbf{WHY OTHER OPTIONS ARE WRONG:}}
		\begin{itemize}
			\item \textcolor{myred}{A:} Responsibility gaps ARE a problem
			\item \textcolor{myred}{C:} Responsibility IS an issue
			\item \textcolor{myred}{D:} Not all employees have clear responsibility
		\end{itemize}
		
		\textcolor{myred}{\textbf{EXAM TRAP:}} GenAI adoption creates \textcolor{myblue}{responsibility gaps} across teams!
	\end{frame}
	
	\begin{frame}{Q95: Misaligned Incentives - Tutorial}
		\fontsize{6.5}{7.5}\selectfont
		\textbf{QUESTION: What misalignment exists in GenAI adoption?}
		
		\vspace{0.1cm}
		\textbf{OPTIONS:}
		\begin{enumerate}
			\item Companies always prioritize safety over speed
			\item Pressure to innovate can outweigh cautionary approaches
			\item Regulation always prevents problems
			\item Incentives are always aligned with safety
		\end{enumerate}
		
		\vspace{0.1cm}
		\textcolor{myblue}{\textbf{ANSWER: B}}
		
		\vspace{0.1cm}
		\textbf{DETAILED EXPLANATION:}
		\begin{itemize}
			\item Pressure to cut costs and improve productivity
			\item Can outweigh safety considerations
			\item Example: Anthropic lobbied to change safety bill
			\item Teams rewarded for cost savings, not risk flagging
			\item Need review structures tied to responsible AI
		\end{itemize}
		
		\vspace{0.1cm}
		\textcolor{myred}{\textbf{WHY OTHER OPTIONS ARE WRONG:}}
		\begin{itemize}
			\item \textcolor{myred}{A:} Safety is NOT always prioritized
			\item \textcolor{myred}{C:} Regulation is not always effective
			\item \textcolor{myred}{D:} Incentives are NOT always aligned
		\end{itemize}
		
		\textcolor{myred}{\textbf{EXAM TRAP:}} Innovation pressure can \textcolor{myblue}{outweigh safety considerations}!
	\end{frame}
	
	\begin{frame}{Q96: Lack of Training Risk - Tutorial}
		\fontsize{6.5}{7.5}\selectfont
		\textbf{QUESTION: What risk arises from lack of GenAI training?}
		
		\vspace{0.1cm}
		\textbf{OPTIONS:}
		\begin{enumerate}
			\item Employees always use GenAI correctly
			\item Employees may use GenAI without understanding acceptable use or data protection obligations
			\item Training is never needed
			\item GenAI always protects data
		\end{enumerate}
		
		\vspace{0.1cm}
		\textcolor{myblue}{\textbf{ANSWER: B}}
		
		\vspace{0.1cm}
		\textbf{DETAILED EXPLANATION:}
		\begin{itemize}
			\item LLM chatbots are intuitive
			\item Employees integrate into workflows without guidance
			\item May paste sensitive data into ChatGPT (4.2% of workers)
			\item Risk to data protection, reputational risk
			\item Companies like JPMorgan banned ChatGPT
		\end{itemize}
		
		\vspace{0.1cm}
		\textcolor{myred}{\textbf{WHY OTHER OPTIONS ARE WRONG:}}
		\begin{itemize}
			\item \textcolor{myred}{A:} Employees may NOT use GenAI correctly
			\item \textcolor{myred}{C:} Training IS needed
			\item \textcolor{myred}{D:} GenAI does NOT always protect data
		\end{itemize}
		
		\textcolor{myred}{\textbf{EXAM TRAP:}} Lack of training leads to \textcolor{myblue}{unintentional data exposure}!
	\end{frame}
	
	\begin{frame}{Q97: Vendor Lock-In - Tutorial}
		\fontsize{6.5}{7.5}\selectfont
		\textbf{QUESTION: What is vendor lock-in in GenAI?}
		
		\vspace{0.1cm}
		\textbf{OPTIONS:}
		\begin{enumerate}
			\item Complete freedom to switch vendors
			\item Dependence on external providers, limiting auditability and leverage
			\item Vendors never change pricing
			\item All providers are identical
		\end{enumerate}
		
		\vspace{0.1cm}
		\textcolor{myblue}{\textbf{ANSWER: B}}
		
		\vspace{0.1cm}
		\textbf{DETAILED EXPLANATION:}
		\begin{itemize}
			\item Most organizations rely on external providers
			\item Few dominant suppliers
			\item Dependencies create power asymmetry
			\item Providers may change pricing/usage without warning
			\item Short-term adoption $\to$ long-term structural lock-in
		\end{itemize}
		
		\vspace{0.1cm}
		\textcolor{myred}{\textbf{WHY OTHER OPTIONS ARE WRONG:}}
		\begin{itemize}
			\item \textcolor{myred}{A:} Switching is often difficult
			\item \textcolor{myred}{C:} Vendors CAN change terms
			\item \textcolor{myred}{D:} Providers differ significantly
		\end{itemize}
		
		\textcolor{myred}{\textbf{EXAM TRAP:}} Vendor lock-in = \textcolor{myblue}{dependence limiting control}!
	\end{frame}
	
	\begin{frame}{Q98: Synthetic Content Risks - Tutorial}
		\fontsize{6.5}{7.5}\selectfont
		\textbf{QUESTION: What risk does synthetic content pose?}
		
		\vspace{0.1cm}
		\textbf{OPTIONS:}
		\begin{enumerate}
			\item Synthetic content is always beneficial
			\item High-quality misleading content can be created at scale, damaging trust and institutions
			\item Synthetic content is easy to identify
			\item Synthetic content never affects elections
		\end{enumerate}
		
		\vspace{0.1cm}
		\textcolor{myblue}{\textbf{ANSWER: B}}
		
		\vspace{0.1cm}
		\textbf{DETAILED EXPLANATION:}
		\begin{itemize}
			\item GenAI lowers cost of creating misleading content
			\item Fake news, health misinformation, political content
			\item Can damage democratic institutions
			\item Can lead to distrust of legitimate sources
			\item EU Code of Practice on Disinformation addresses risks
		\end{itemize}
		
		\vspace{0.1cm}
		\textcolor{myred}{\textbf{WHY OTHER OPTIONS ARE WRONG:}}
		\begin{itemize}
			\item \textcolor{myred}{A:} Synthetic content CAN be harmful
			\item \textcolor{myred}{C:} Synthetic content is often difficult to identify
			\item \textcolor{myred}{D:} Synthetic content CAN affect elections
		\end{itemize}
		
		\textcolor{myred}{\textbf{EXAM TRAP:}} Synthetic content can \textcolor{myblue}{damage trust and institutions}!
	\end{frame}
	
	\begin{frame}{Q99: Deepfakes Risks - Tutorial}
		\fontsize{6.5}{7.5}\selectfont
		\textbf{QUESTION: What risk do deepfakes pose?}
		
		\vspace{0.1cm}
		\textbf{OPTIONS:}
		\begin{enumerate}
			\item Deepfakes are always harmless
			\item Audio or video deepfakes can convincingly mimic individuals, causing fraud, defamation, or political manipulation
			\item Deepfakes are easy to detect
			\item Deepfakes never cause harm
		\end{enumerate}
		
		\vspace{0.1cm}
		\textcolor{myblue}{\textbf{ANSWER: B}}
		
		\vspace{0.1cm}
		\textbf{DETAILED EXPLANATION:}
		\begin{itemize}
			\item Convincingly mimic public figures or private individuals
			\item Individual and societal harm
			\item Fraud: CEO persuaded to transfer €220,000
			\item Defamation, political manipulation
			\item Even low-effort impersonations can go viral
		\end{itemize}
		
		\vspace{0.1cm}
		\textcolor{myred}{\textbf{WHY OTHER OPTIONS ARE WRONG:}}
		\begin{itemize}
			\item \textcolor{myred}{A:} Deepfakes CAN cause harm
			\item \textcolor{myred}{C:} Deepfakes are often difficult to detect
			\item \textcolor{myred}{D:} Deepfakes CAN cause harm
		\end{itemize}
		
		\textcolor{myred}{\textbf{EXAM TRAP:}} Deepfakes can \textcolor{myblue}{cause fraud, defamation, and manipulation}!
	\end{frame}
	
	\begin{frame}{Q100: GenAI Risk Management Summary - Tutorial}
		\fontsize{6.5}{7.5}\selectfont
		\textbf{QUESTION: What is the most comprehensive summary of GenAI risk management?}
		
		\vspace{0.1cm}
		\textbf{OPTIONS:}
		\begin{enumerate}
			\item GenAI risks are purely technical
			\item GenAI risks include functional, data-related, organizational, strategic, and societal risks, requiring integrated oversight and governance
			\item GenAI risks are not significant
			\item Only hallucinations matter
		\end{enumerate}
		
		\vspace{0.1cm}
		\textcolor{myblue}{\textbf{ANSWER: B}}
		
		\vspace{0.1cm}
		\textbf{DETAILED EXPLANATION:}
		\begin{itemize}
			\item \textbf{Functional:} Hallucinations, unpredictability, adversarial inputs
			\item \textbf{Data-related:} Privacy violations, copyright issues
			\item \textbf{Organizational:} Responsibility gaps, misaligned incentives, lack of training
			\item \textbf{Strategic:} Vendor lock-in, regulatory unpredictability
			\item \textbf{Societal:} Disinformation, deepfakes, manipulation
		\end{itemize}
		
		\vspace{0.1cm}
		\textcolor{myred}{\textbf{WHY OTHER OPTIONS ARE WRONG:}}
		\begin{itemize}
			\item \textcolor{myred}{A:} Risks extend beyond technical
			\item \textcolor{myred}{C:} GenAI risks ARE significant
			\item \textcolor{myred}{D:} Multiple risk categories exist
		\end{itemize}
		
		\textcolor{myred}{\textbf{EXAM SUCCESS:}} Understand \textcolor{myblue}{all risk categories} and their interconnections!
	\end{frame}
	
\end{document}